% Options for packages loaded elsewhere
\PassOptionsToPackage{unicode}{hyperref}
\PassOptionsToPackage{hyphens}{url}
\documentclass[
]{article}
\usepackage{xcolor}
\usepackage[margin=1in]{geometry}
\usepackage{amsmath,amssymb}
\setcounter{secnumdepth}{-\maxdimen} % remove section numbering
\usepackage{iftex}
\ifPDFTeX
  \usepackage[T1]{fontenc}
  \usepackage[utf8]{inputenc}
  \usepackage{textcomp} % provide euro and other symbols
\else % if luatex or xetex
  \usepackage{unicode-math} % this also loads fontspec
  \defaultfontfeatures{Scale=MatchLowercase}
  \defaultfontfeatures[\rmfamily]{Ligatures=TeX,Scale=1}
\fi
\usepackage{lmodern}
\ifPDFTeX\else
  % xetex/luatex font selection
\fi
% Use upquote if available, for straight quotes in verbatim environments
\IfFileExists{upquote.sty}{\usepackage{upquote}}{}
\IfFileExists{microtype.sty}{% use microtype if available
  \usepackage[]{microtype}
  \UseMicrotypeSet[protrusion]{basicmath} % disable protrusion for tt fonts
}{}
\makeatletter
\@ifundefined{KOMAClassName}{% if non-KOMA class
  \IfFileExists{parskip.sty}{%
    \usepackage{parskip}
  }{% else
    \setlength{\parindent}{0pt}
    \setlength{\parskip}{6pt plus 2pt minus 1pt}}
}{% if KOMA class
  \KOMAoptions{parskip=half}}
\makeatother
\usepackage{longtable,booktabs,array}
\newcounter{none} % for unnumbered tables
\usepackage{calc} % for calculating minipage widths
% Correct order of tables after \paragraph or \subparagraph
\usepackage{etoolbox}
\makeatletter
\patchcmd\longtable{\par}{\if@noskipsec\mbox{}\fi\par}{}{}
\makeatother
% Allow footnotes in longtable head/foot
\IfFileExists{footnotehyper.sty}{\usepackage{footnotehyper}}{\usepackage{footnote}}
\makesavenoteenv{longtable}
\usepackage{graphicx}
\makeatletter
\newsavebox\pandoc@box
\newcommand*\pandocbounded[1]{% scales image to fit in text height/width
  \sbox\pandoc@box{#1}%
  \Gscale@div\@tempa{\textheight}{\dimexpr\ht\pandoc@box+\dp\pandoc@box\relax}%
  \Gscale@div\@tempb{\linewidth}{\wd\pandoc@box}%
  \ifdim\@tempb\p@<\@tempa\p@\let\@tempa\@tempb\fi% select the smaller of both
  \ifdim\@tempa\p@<\p@\scalebox{\@tempa}{\usebox\pandoc@box}%
  \else\usebox{\pandoc@box}%
  \fi%
}
% Set default figure placement to htbp
\def\fps@figure{htbp}
\makeatother
\setlength{\emergencystretch}{3em} % prevent overfull lines
\providecommand{\tightlist}{%
  \setlength{\itemsep}{0pt}\setlength{\parskip}{0pt}}
\usepackage[]{natbib}
\bibliographystyle{plainnat}
\usepackage{longtable,booktabs}\DeclareUnicodeCharacter{2260}{\ensuremath{\neq}}\DeclareUnicodeCharacter{2265}{\ensuremath{\geq}}\DeclareUnicodeCharacter{2264}{\ensuremath{\leq}}\DeclareUnicodeCharacter{2192}{\ensuremath{\rightarrow}}\DeclareUnicodeCharacter{2282}{\ensuremath{\subset}}\DeclareUnicodeCharacter{2248}{\ensuremath{\approx}}\DeclareUnicodeCharacter{00D7}{\ensuremath{\times}}\DeclareUnicodeCharacter{2211}{\ensuremath{\sum}}\DeclareUnicodeCharacter{221E}{\ensuremath{\infty}}\DeclareUnicodeCharacter{03B5}{\ensuremath{\varepsilon}}\DeclareUnicodeCharacter{03C4}{\ensuremath{\tau}}\DeclareUnicodeCharacter{03BB}{\ensuremath{\lambda}}\DeclareUnicodeCharacter{03A9}{\ensuremath{\Omega}}\DeclareUnicodeCharacter{03B8}{\ensuremath{\theta}}\DeclareUnicodeCharacter{2208}{\ensuremath{\in}}\DeclareUnicodeCharacter{2026}{\ldots}\DeclareUnicodeCharacter{2020}{\dag}\DeclareUnicodeCharacter{2021}{\ddag}\DeclareUnicodeCharacter{00A7}{\S}\DeclareUnicodeCharacter{03BA}{\ensuremath{\kappa}}\DeclareUnicodeCharacter{03B1}{\ensuremath{\alpha}}\DeclareUnicodeCharacter{03B2}{\ensuremath{\beta}}\DeclareUnicodeCharacter{03B4}{\ensuremath{\delta}}\DeclareUnicodeCharacter{03B3}{\ensuremath{\gamma}}\DeclareUnicodeCharacter{2295}{\ensuremath{\oplus}}\DeclareUnicodeCharacter{2218}{\ensuremath{\circ}}\DeclareUnicodeCharacter{2212}{-}\DeclareUnicodeCharacter{2019}{\textquoteright}\DeclareUnicodeCharacter{2018}{\textquoteleft}\DeclareUnicodeCharacter{201C}{\textquotedblleft}\DeclareUnicodeCharacter{201D}{\textquotedblright}\DeclareUnicodeCharacter{2070}{\ensuremath{^0}}\DeclareUnicodeCharacter{00B1}{\ensuremath{\pm}}\DeclareUnicodeCharacter{2264}{\ensuremath{\leq}}
\usepackage{bookmark}
\IfFileExists{xurl.sty}{\usepackage{xurl}}{} % add URL line breaks if available
\urlstyle{same}
\hypersetup{
  pdftitle={Supplementary Material for: Frontier and Localhost: A Survey of How Production AI Learns Outside the Weights},
  hidelinks,
  pdfcreator={LaTeX via pandoc}}

\title{Supplementary Material for: Frontier and Localhost: A Survey of
How Production AI Learns Outside the Weights}
\author{}
\date{October 2026}

\begin{document}
\maketitle

This document accompanies the main article. S1 is the released score
table; S2--S5 carry the verification material of the survey (data
appendix, mechanism stack, contrast set and matrix membership, and the
per-system walkthrough of the convergent research families); S6 is the
audit of every absolute claim of the earlier preprint against the
October-2026 literature, with the rubric-reward comparison that Section
IV of the main text summarizes. Section numbers in the form ``Section
N'' refer to the main article.

\section{S1. Full score table}\label{s1.-full-score-table}

Every scored (substrate, system) row of the corpus: 130 systems scored
by May 2026 and 126 added by the October 2026 sweep. Scores follow the
0--5 rubric of the main text (Table V-1); Loop 1 / Loop 2 gate types and
signal classes are recorded for the October rows (and for May rows
re-audited in Section VII). Substrate codes: INT integrated, S1
instructions, S2 skills, S3 memory, S4 tools, S5 orchestration, S6
governance. The CSV files with per-row justifications accompany this
document.

\begingroup\scriptsize

{\def\LTcaptype{none} % do not increment counter
\begin{longtable}[]{@{}
  >{\raggedright\arraybackslash}p{(\linewidth - 16\tabcolsep) * \real{0.0600}}
  >{\raggedright\arraybackslash}p{(\linewidth - 16\tabcolsep) * \real{0.3600}}
  >{\raggedright\arraybackslash}p{(\linewidth - 16\tabcolsep) * \real{0.0400}}
  >{\raggedright\arraybackslash}p{(\linewidth - 16\tabcolsep) * \real{0.0400}}
  >{\raggedright\arraybackslash}p{(\linewidth - 16\tabcolsep) * \real{0.1800}}
  >{\raggedright\arraybackslash}p{(\linewidth - 16\tabcolsep) * \real{0.1000}}
  >{\raggedright\arraybackslash}p{(\linewidth - 16\tabcolsep) * \real{0.1000}}
  >{\raggedright\arraybackslash}p{(\linewidth - 16\tabcolsep) * \real{0.0400}}
  >{\raggedright\arraybackslash}p{(\linewidth - 16\tabcolsep) * \real{0.0800}}@{}}
\toprule\noalign{}
\begin{minipage}[b]{\linewidth}\raggedright
Sub.
\end{minipage} & \begin{minipage}[b]{\linewidth}\raggedright
System
\end{minipage} & \begin{minipage}[b]{\linewidth}\raggedright
Score
\end{minipage} & \begin{minipage}[b]{\linewidth}\raggedright
Year
\end{minipage} & \begin{minipage}[b]{\linewidth}\raggedright
Source
\end{minipage} & \begin{minipage}[b]{\linewidth}\raggedright
Loop 1
\end{minipage} & \begin{minipage}[b]{\linewidth}\raggedright
Loop 2
\end{minipage} & \begin{minipage}[b]{\linewidth}\raggedright
Class
\end{minipage} & \begin{minipage}[b]{\linewidth}\raggedright
Added
\end{minipage} \\
\midrule\noalign{}
\endhead
\bottomrule\noalign{}
\endlastfoot
INT & ADAS / Meta Agent Search & 4 & 2024 & arXiv:2408.08435 & & & & May
26 \\
INT & Agent S / S2 & 4.5 & 2024 & arXiv:2410.08164 & & & & May 26 \\
INT & Agent Workflow Memory (AWM) & 3.5 & 2024 & arXiv:2409.07429 & & &
& May 26 \\
INT & AlphaEvolve & 4 & 2025 & arXiv:2506.13131 & & & & May 26 \\
INT & AutoAgent & 5 & 2026 & arXiv:2603.09716 & & & & May 26 \\
INT & AutoManual & 3.5 & 2024 & NeurIPS 2024 & & & & May 26 \\
INT & AutoSaddler & 4 & 2026 & arXiv:2608.23041 & Auto & None & C2 & Oct
26 \\
INT & CaptainAgent & 2.5 & 2024 & arXiv:2405.19425 & & & & May 26 \\
INT & CASCADE & 4.5 & 2025 & arXiv:2512.23880 & & & & May 26 \\
INT & ChatDev & 3 & 2023 & arXiv:2307.07924 & & & & May 26 \\
INT & Co-Harness & 4 & 2026 & arXiv:2607.22688 & Auto & None & C2 & Oct
26 \\
INT & CoEvoSkills & 3.5 & 2026 & arXiv:2604.01687 & & & & May 26 \\
INT & Cognition Devin & 3 & 2024 & cognition.ai & & & & May 26 \\
INT & Continual Harness & 4 & 2026 & arXiv:2605.09998 & None & None & C4
& Oct 26 \\
INT & Cradle & 4 & 2024 & arXiv:2403.03186 & & & & May 26 \\
INT & Darwin Gödel Machine & 4 & 2025 & arXiv:2505.22954 & & & & May
26 \\
INT & Double Ratchet & 4 & 2026 & arXiv:2607.12790 & Auto & None & C3 an
& Oct 26 \\
INT & EvolveNet & 4 & 2026 & arXiv:2608.04968 & Auto & Auto & C1,C2 &
Oct 26 \\
INT & EvolveR & 4 & 2025 & arXiv:2510.16079 & & & & May 26 \\
INT & ExpeL & 3.5 & 2023 & arXiv:2308.10144 & & & & May 26 \\
INT & Generative Agents & 4 & 2023 & arXiv:2304.03442 & & & & May 26 \\
INT & GenericAgent & 4 & 2026 & arXiv:2604.17091 & Auto & None & C2 &
Oct 26 \\
INT & GITM & 3 & 2023 & arXiv:2305.17144 & & & & May 26 \\
INT & HarnessBank & 4 & 2026 & arXiv:2607.13683 & Auto & None & C2 & Oct
26 \\
INT & HarnessX & 4 & 2026 & arXiv:2606.14249 & Auto & None & C2 & Oct
26 \\
INT & Hippocratic AI Polaris & 3.5 & 2024 & hippocraticai.com & & & &
May 26 \\
INT & JARVIS-1 & 3.5 & 2023 & arXiv:2311.05997 & & & & May 26 \\
INT & JIT-Agent & 4 & 2026 & arXiv:2608.25593 & Auto & None & C2 & Oct
26 \\
INT & Life-Harness & 4 & 2026 & arXiv:2605.22166 & Auto & None & C2 &
Oct 26 \\
INT & Mendel Gödel Machine (MGM) & 4 & 2026 & arXiv:2608.07645 & Auto &
None & C1 & Oct 26 \\
INT & MetaGPT & 3.5 & 2023 & arXiv:2308.00352 & & & & May 26 \\
INT & Mobile-Agent-E & 4.5 & 2025 & arXiv:2501.11733 & & & & May 26 \\
INT & ModularRSI & 4 & 2026 & arXiv:2609.14857 & Auto & None & C2 & Oct
26 \\
INT & NanoResearch & 5 & 2026 & arXiv:2605.10813 & & & & May 26 \\
INT & OpenCore & 4 & 2025 & github.com/sibmike/opencore & & & & May
26 \\
INT & OpenHands / CodeAct & 3 & 2024 & arXiv:2407.16741 & & & & May
26 \\
INT & OS-Copilot / FRIDAY & 4 & 2024 & arXiv:2402.07456 & & & & May
26 \\
INT & Plan4MC & 2.5 & 2023 & GitHub/arXiv & & & & May 26 \\
INT & SAGE & 4 & 2025 & arXiv:2512.17102 & & & & May 26 \\
INT & Sakana AI Scientist & 3 & 2024 & arXiv:2408.06292 & & & & May
26 \\
INT & ScholarEvolve & 4 & 2026 & arXiv:2609.40169 & Auto & None & C2 &
Oct 26 \\
INT & SIA & 4 & 2026 & arXiv:2605.27276 & Auto & None & C2 & Oct 26 \\
INT & Sierra Agent OS & 4 & 2024 & sierra.ai & & & & May 26 \\
INT & Sierra Agent OS 2.0 & 3 & 2025 & sierra.ai/blog/agent-os-2-0 & & &
& May 26 \\
INT & SkillForge & 4 & 2026 & arXiv:2604.08618 & & & & May 26 \\
INT & SkillMAS & 4 & 2026 & arXiv:2605.09341 & Auto & None & C2 & Oct
26 \\
INT & SkillOS & 3 & 2026 & SemanticScholar:2653e4c7 & & & & May 26 \\
INT & SkillRL & 5 & 2026 & arXiv:2602.08234 & & & & May 26 \\
INT & TTHE & 4 & 2026 & arXiv:2607.08124 & Auto & None & C3 & Oct 26 \\
INT & VACE & 4 & 2026 & arXiv:2609.37105 & Auto & None & C2 & Oct 26 \\
INT & Voyager & 4 & 2023 & arXiv:2305.16291 & & & & May 26 \\
INT & Workspace Optimization (DreamTeam) & 3 & 2026 & arXiv:2605.09650 &
Auto & None & C1 & Oct 26 \\
S1 & AGENTS.md & 2 & 2026 & arXiv:2601.20404 & & & & May 26 \\
S1 & Aider Conventions & 2 & 2024 & aider.chat & & & & May 26 \\
S1 & Claude Code CLAUDE.md & 4 & 2025 & docs.anthropic.com & & & & May
26 \\
S1 & Cody / Sourcegraph & 2 & 2024 & sourcegraph.com & & & & May 26 \\
S1 & Continue.dev Rules & 3 & 2024 & docs.continue.dev & & & & May 26 \\
S1 & Cursor Rules & 3 & 2024 & docs.cursor.com/rules & & & & May 26 \\
S1 & GitHub Copilot Custom Instructions & 3 & 2024 & docs.github.com & &
& & May 26 \\
S1 & GRACE & 4 & 2026 & arXiv:2607.09175 & Auto & None & C2 & Oct 26 \\
S1 & Lovable Knowledge & 2 & 2024 & docs.lovable.dev & & & & May 26 \\
S1 & Meta Context Engineering (MCE) & 4 & 2026 & arXiv:2601.21557 & Auto
& None & C2 & Oct 26 \\
S1 & Replit replit.md & 4 & 2025 & docs.replit.com & & & & May 26 \\
S1 & Spotify conversational-recommendation loop & 4 & 2026 &
arXiv:2609.30297 & Auto & Human & C3 & Oct 26 \\
S1 & VS Code Copilot Instructions & 3 & 2024 & code.visualstudio.com & &
& & May 26 \\
S1 & Windsurf Cascade Rules & 3 & 2024 & docs.windsurf.com & & & & May
26 \\
S1 & Zed AI Rules & 3 & 2024 & zed.dev/docs & & & & May 26 \\
S2 & AgentVerse & 1 & 2023 & arXiv:2308.10848 & & & & May 26 \\
S2 & Anthropic Claude Agent Skills & 3 & 2025 &
anthropic.com/engineering & & & & May 26 \\
S2 & Anthropic/skills GitHub repo & 2 & 2025 &
github.com/anthropics/skills & & & & May 26 \\
S2 & ASSAY & 4 & 2026 & arXiv:2606.15390 & Auto & None & C2 & Oct 26 \\
S2 & AutoSkill & 3 & 2026 & arXiv:2603.01145 & Auto & None & C3 & Oct
26 \\
S2 & Buildrix & 3 & 2026 & arXiv:2606.25139 & Human & Human & C3 & Oct
26 \\
S2 & ChatGPT Plugins (legacy) & 1 & 2023 & openai.com & & & & May 26 \\
S2 & Claude Code Skills & 3 & 2025 & docs.anthropic.com & & & & May
26 \\
S2 & COSPLAY & 5 & 2026 & arXiv:2604.20987 & & & & May 26 \\
S2 & CrewAI Skills & 2 & 2025 & docs.crewai.com & & & & May 26 \\
S2 & HiSME & 4 & 2026 & arXiv:2605.28390 & Auto & None & C1,C3 & Oct
26 \\
S2 & HuggingFace Agent Skills & 2 & 2025 & huggingface.co & & & & May
26 \\
S2 & LangChain Skills & 2 & 2025 & langchain.com & & & & May 26 \\
S2 & Library Drift curator & 4 & 2026 & arXiv:2605.19576 & Auto & None &
C1 & Oct 26 \\
S2 & Meta Context Engineering (MCE) & 4 & 2026 & arXiv:2601.21557 & Auto
& None & C2 & Oct 26 \\
S2 & MetaGPT (SOP-as-skill) & 3 & 2023 & arXiv:2308.00352 & & & & May
26 \\
S2 & MetaSkill-Evolve & 4 & 2026 & arXiv:2607.05297 & Auto & None & C2 &
Oct 26 \\
S2 & Microsoft Semantic Kernel Plugins & 2 & 2024 & learn.microsoft.com
& & & & May 26 \\
S2 & MUSE-Autoskill & 4 & 2026 & arXiv:2605.27366 & Auto & None & C1,C2
& Oct 26 \\
S2 & NanoResearch Skill Bank & 4 & 2026 & arXiv:2605.10813 & & & & May
26 \\
S2 & OpenAI Custom GPTs / Actions & 2 & 2023 & platform.openai.com & & &
& May 26 \\
S2 & Personalized developer skills (Huang-Du-Lan) & 3 & 2026 &
arXiv:2608.10319 & Auto & None & C3 & Oct 26 \\
S2 & RASO & 4 & 2026 & arXiv:2609.38024 & Auto & Auto & C2 & Oct 26 \\
S2 & SAGE & 5 & 2025 & arXiv:2512.17102 & & & & May 26 \\
S2 & SCAFFOLD & 4 & 2026 & arXiv:2609.05511 & Auto & None & C2 & Oct
26 \\
S2 & SE-GoS & 4 & 2026 & arXiv:2609.08228 & Auto & None & C2 & Oct 26 \\
S2 & SkeMex & 4 & 2026 & arXiv:2606.09365 & Auto & None & C2 & Oct 26 \\
S2 & Skill Issues (Bauplan) & 4 & 2026 & arXiv:2606.01185 & Auto & None
& C1,C3 & Oct 26 \\
S2 & SKILL.md smells (Hong-Imani-Ahmed) & 1 & 2026 & arXiv:2607.01456 &
None & None & - & Oct 26 \\
S2 & SkillAudit & 4 & 2026 & arXiv:2606.14239 & Auto & None & C3,C1 &
Oct 26 \\
S2 & SkillBoost & 4 & 2026 & arXiv:2607.26643 & Auto & None & C1 & Oct
26 \\
S2 & SkillChain & 4 & 2026 & arXiv:2606.12984 & Auto-\textgreater Huma &
None & C3,C2 & Oct 26 \\
S2 & SkillComposer & 4 & 2026 & arXiv:2606.06079 & Auto & None & C1,C2 &
Oct 26 \\
S2 & SkillEvoReg & 4 & 2026 & arXiv:2609.30861 & Auto & None & C1,C2 &
Oct 26 \\
S2 & SkillFoundry & 3 & 2026 & arXiv:2604.03964 & Auto & None & C1 & Oct
26 \\
S2 & SkillGen & 4 & 2026 & arXiv:2605.10999 & Auto & None & C2 & Oct
26 \\
S2 & SkillGLoW & 4 & 2026 & arXiv:2609.02217 & Auto & None & C2 & Oct
26 \\
S2 & SkillLift & 4 & 2026 & arXiv:2609.15396 & Auto & None & C3 in & Oct
26 \\
S2 & SkillOps & 4 & 2026 & arXiv:2605.13716 & Auto & None & C2 & Oct
26 \\
S2 & SkillSmith & 4 & 2026 & arXiv:2606.01314 & Auto & None & C1,C2 &
Oct 26 \\
S2 & SkillSpec (consensus-gated) & 4 & 2026 & arXiv:2610.00704 & Auto &
None & C1,C2 & Oct 26 \\
S2 & SuperAGI Toolkits & 1 & 2024 & github.com/TransformerOptimu & & & &
May 26 \\
S2 & TRUSS & 3 & 2026 & arXiv:2608.17588 & Auto & None & C1,C2 & Oct
26 \\
S2 & VALVE & 4 & 2026 & arXiv:2609.32990 & Auto & None & C1 & Oct 26 \\
S2 & VCE-Skill & 4 & 2026 & arXiv:2608.16544 & Auto & None & C1,C2 & Oct
26 \\
S2 & Voyager & 4 & 2023 & arXiv:2305.16291 & & & & May 26 \\
S2 & Web2BigTable & 3 & 2026 & arXiv:2604.27221 & Auto & None & C1,C2 &
Oct 26 \\
S2 & Who Maintains Agent Skills? (empirical) & 2 & 2026 &
arXiv:2609.05677 & Human & None & C3 & Oct 26 \\
S2 & WikiSkill & 4 & 2026 & arXiv:2608.27454 & Auto & None & C2 & Oct
26 \\
S2 & XRepoSkill & 4 & 2026 & arXiv:2609.36807 & Auto & Auto-parti & C1 &
Oct 26 \\
S3 & A-MEM & 3 & 2025 & NeurIPS 2025 & & & & May 26 \\
S3 & Agent-native memory system evaluation (Zhou et a & 2 & 2026 &
arXiv:2606.24775 & n/a & n/a & n/a & Oct 26 \\
S3 & AIM privacy-aware interoperable memory (Johnson & 3 & 2026 &
arXiv:2609.12320 & Auto & None & C3 & Oct 26 \\
S3 & Beyond the Harness (enterprise Text-to-SQL cards & 3 & 2026 &
arXiv:2608.22830 & Auto & None & C2 & Oct 26 \\
S3 & ChatGPT Memory & 3 & 2024 & openai.com & & & & May 26 \\
S3 & Claude Memory & 3 & 2025 & anthropic.com & & & & May 26 \\
S3 & Cognee & 3 & 2024 & github.com/topoteretes/cogne & & & & May 26 \\
S3 & Compiled Memory (Rhodes \& Kang) & 4 & 2026 & arXiv:2603.15666 &
Auto-gated & None & C1 & Oct 26 \\
S3 & Control-plane placement / ForgetEval (Yang) & 2 & 2026 &
arXiv:2606.15903 & n/a & n/a & n/a & Oct 26 \\
S3 & FRESH failure-aware heterogeneous graph memory ( & 4 & 2026 &
arXiv:2609.28003 & Auto & None & C2 & Oct 26 \\
S3 & GateMem multi-principal memory-governance benchm & 2 & 2026 &
arXiv:2606.18829 & n/a & n/a & n/a & Oct 26 \\
S3 & Gemini Code Assist Memory & 4 & 2025 & cloud.google.com & & & & May
26 \\
S3 & Generative Agents & 3 & 2023 & arXiv:2304.03442 & & & & May 26 \\
S3 & Governed Collaborative Memory (Cuadros et al.) & 5 & 2026 &
arXiv:2605.04264 & & & & May 26 \\
S3 & HippoRAG & 2 & 2024 & arXiv & & & & May 26 \\
S3 & Honest Lying mitigation / programmatic failure s & 4 & 2026 &
arXiv:2605.29463 & Auto & None & C1/C2 & Oct 26 \\
S3 & KV-cache rollback consistency (Zhang \& Yang) & 2 & 2026 &
arXiv:2608.15939 & n/a & n/a & n/a & Oct 26 \\
S3 & LangChain Primitives (legacy) & 1 & 2023 & langchain.com & & & &
May 26 \\
S3 & LangGraph Long-Term Memory & 3 & 2024 & langchain.com & & & & May
26 \\
S3 & LangGraph Store / LangMem (2026 refresh) & 3 & 2026 &
docs.langchain.com/oss/pytho & None & None & n/a & Oct 26 \\
S3 & Learning on the Job & 4 & 2026 & arXiv:2607.22157 & Auto+Human &
None & C2,C3 & Oct 26 \\
S3 & Letta MemFS + Context Repositories + Dreaming (2 & 4 & 2026 &
docs.letta.com/letta-code/me & Auto & Auto & C3 & Oct 26 \\
S3 & Mem0 & 3 & 2025 & arXiv & & & & May 26 \\
S3 & Mem0 Platform (2026 refresh: history, feedback, & 3 & 2026 &
docs.mem0.ai/api-reference/m & Auto & None & C3 & Oct 26 \\
S3 & MemAudit post-hoc poisoned-memory audit (Tan et & 3 & 2026 &
arXiv:2605.23723 & Auto & None & C2 & Oct 26 \\
S3 & MemClaw / Governed Shared Memory (Margalit et al & 3 & 2026 &
arXiv:2606.24535 & Auto & None & C3 & Oct 26 \\
S3 & MemGPT / Letta & 3 & 2023 & arXiv:2310.08560 & & & & May 26 \\
S3 & MemGuard persisted verifier signals (Wang et al. & 4 & 2026 &
arXiv:2608.21867 & Auto-gated & None & C3 & Oct 26 \\
S3 & MemoryBank & 3 & 2023 & Zhong et al. & & & & May 26 \\
S3 & Mi-Memory lifecycle framework (Liu et al.; Xiaom & 3 & 2026 &
arXiv:2607.18975 & Auto & None & C3 & Oct 26 \\
S3 & MIRIX & 3 & 2025 & arXiv & & & & May 26 \\
S3 & MPBench memory-poisoning study (Dash et al.) & 2 & 2026 &
arXiv:2606.04329 & n/a & n/a & n/a & Oct 26 \\
S3 & NanoResearch Memory Module & 3 & 2026 & arXiv:2605.10813 & & & &
May 26 \\
S3 & PPMF provenance-preserving memory firewall (Xu e & 3 & 2026 &
arXiv:2607.29167 & Auto & None & C1 & Oct 26 \\
S3 & REALM retrieval-driven memory reconsolidation (S & 4 & 2026 &
arXiv:2609.16053 & Auto & None & C2/C3 & Oct 26 \\
S3 & Recuris recursive experiential-working memory (Y & 4 & 2026 &
arXiv:2608.24876 & Auto-gated & None & C2 & Oct 26 \\
S3 & Reflexion & 4 & 2023 & arXiv:2303.11366 & & & & May 26 \\
S3 & Revocation guard / Revoked but Still Authoritati & 3 & 2026 &
arXiv:2609.08258 & Human & None & C1 & Oct 26 \\
S3 & Safe to Resume? checkpoint-rollback continuity ( & 2 & 2026 &
arXiv:2608.29381 & n/a & n/a & n/a & Oct 26 \\
S3 & ShareMem & 4 & 2026 & arXiv:2609.32511 & Auto & Auto-parti & C3 &
Oct 26 \\
S3 & SIx Harness / Memory as Infrastructure (Helwig) & 3 & 2026 &
arXiv:2609.05510 & Auto & None & C2 & Oct 26 \\
S3 & SOAR Cognitive Architecture & 4 & classical & Newell/Laird & & & &
May 26 \\
S3 & SSGM stability-and-safety governed memory (Lam e & 2 & 2026 &
arXiv:2603.11768 & None & None & n/a & Oct 26 \\
S3 & Stale constraints in inherited memory (Nakayashi & 2 & 2026 &
arXiv:2608.25553 & n/a & n/a & n/a & Oct 26 \\
S3 & TMA-NM origin-bound memory authority (Louck) & 3 & 2026 &
arXiv:2606.24322 & Auto & None & C1 & Oct 26 \\
S3 & TOKI bitemporal operator algebra (Wang) & 3 & 2026 &
arXiv:2606.06240 & Auto or Hu & None & C3 & Oct 26 \\
S3 & Trust-scored memory sanitization (Devarangadi Su & 2 & 2026 &
arXiv:2601.05504 & Auto gate & None & n/a & Oct 26 \\
S3 & Typed provenance + assertion guardrails / Stored & 3 & 2026 &
arXiv:2609.02127 & Auto & None & C1 & Oct 26 \\
S3 & Utility Under Attack (Karunanidhi) & 2 & 2026 & arXiv:2608.21230 &
n/a & n/a & n/a & Oct 26 \\
S3 & VerificAgent human-verified frozen memory (Nguye & 3 & 2025 &
arXiv:2506.02539 & Human & None & C3 & Oct 26 \\
S3 & VerMem verifiable memory management (Sun et al.) & 4 & 2026 &
arXiv:2608.03137 & Auto-gated & None & C2/C3 & Oct 26 \\
S3 & Voyager Skill Library & 4 & 2023 & arXiv:2305.16291 & & & & May
26 \\
S3 & Zep / Graphiti & 3 & 2025 & arXiv & & & & May 26 \\
S3 & Zep / Graphiti bitemporal facts (2026 refresh) & 3 & 2026 &
help.getzep.com/facts; help. & Auto & None & C3 & Oct 26 \\
S4 & Anthropic Tool Use API & 4 & 2024 & docs.anthropic.com & & & & May
26 \\
S4 & Browser Tools (Playwright/Computer Use) & 3 & 2024 & anthropic.com
& & & & May 26 \\
S4 & Claude Code Tool API & 4 & 2024 & docs.anthropic.com & & & & May
26 \\
S4 & Code Interpreter & 3 & 2023 & openai.com & & & & May 26 \\
S4 & Cursor MCP Integration & 4 & 2024 & docs.cursor.com & & & & May
26 \\
S4 & Devin / SWE-agent Tool Harness & 4 & 2024 & cognition.ai & & & &
May 26 \\
S4 & Gemini Function Calling & 3 & 2024 & ai.google.dev & & & & May
26 \\
S4 & GPT Actions & 3 & 2023 & platform.openai.com & & & & May 26 \\
S4 & Harvey (Legal Vertical) & 5 & 2024 & harvey.ai & & & & May 26 \\
S4 & Hippocratic AI & 5 & 2024 & hippocraticai.com & & & & May 26 \\
S4 & LangChain Tools Abstraction & 3 & 2023 & langchain.com & & & & May
26 \\
S4 & LangGraph Tool Nodes & 4 & 2024 & langchain.com & & & & May 26 \\
S4 & LlamaIndex Tool Ecosystem & 4 & 2023 & llamaindex.ai & & & & May
26 \\
S4 & Model Context Protocol (MCP) & 5 & 2024 & modelcontextprotocol.io &
& & & May 26 \\
S4 & OpenAI Assistants API & 3 & 2023 & platform.openai.com & & & & May
26 \\
S4 & OpenAI Function Calling & 3 & 2023 & platform.openai.com & & & &
May 26 \\
S4 & Sierra & 4 & 2024 & sierra.ai & & & & May 26 \\
S4 & SkillSmith & 4 & 2026 & arXiv:2606.01314 & Auto & None & C1 & Oct
26 \\
S4 & Toolformer & 2 & 2023 & arXiv:2302.04761 & & & & May 26 \\
S4 & ToolLLM / ToolBench & 3 & 2023 & arXiv:2307.16789 & & & & May 26 \\
S5 & AgentConductor & 0 & 2026 & arXiv:2602.17100 & None & None & C2 &
Oct 26 \\
S5 & AgentGrad & 4 & 2026 & arXiv:2609.08572 & Auto & None & C2 & Oct
26 \\
S5 & AgentVerse & 2 & 2023 & arXiv:2308.10848 & & & & May 26 \\
S5 & AutoGen / MagenticOne & 3 & 2024 & arXiv:2308.08155 & & & & May
26 \\
S5 & AWS Bedrock Multi-Agent & 2 & 2024 & aws.amazon.com & & & & May
26 \\
S5 & Bilevel Autoresearch & 4 & 2026 & arXiv:2603.23420 & Auto & None &
C1 & Oct 26 \\
S5 & CAMEL & 1 & 2023 & arXiv & & & & May 26 \\
S5 & ChatDev & 2 & 2023 & arXiv:2307.07924 & & & & May 26 \\
S5 & ContextEvo & 4 & 2026 & arXiv:2609.34649 & Auto & None & C2 & Oct
26 \\
S5 & Control System for Frozen Agents (Paul) & 4 & 2026 &
arXiv:2607.25415 & Auto & None & C2 & Oct 26 \\
S5 & CrewAI & 3 & 2025 & docs.crewai.com & & & & May 26 \\
S5 & Evo2Team & 4 & 2026 & arXiv:2609.34135 & Auto & Auto & C2 & Oct
26 \\
S5 & Growing Harness & 4 & 2026 & arXiv:2609.26760 & Auto & None & C2 &
Oct 26 \\
S5 & HARBOR & 4 & 2026 & arXiv:2604.20938 & Auto & None & C2 & Oct 26 \\
S5 & HarnessForge & 4 & 2026 & arXiv:2606.01779 & Auto & None & C2 & Oct
26 \\
S5 & Inspect AI & 2 & 2024 & ukgovernmentbeis.github.io & & & & May
26 \\
S5 & Institutional MAS architectures (Fei et al.) & 1 & 2026 &
arXiv:2604.27691 & None & None & C2 & Oct 26 \\
S5 & LangGraph & 3 & 2024 & langchain.com & & & & May 26 \\
S5 & Letta & 3 & 2025 & letta.com & & & & May 26 \\
S5 & LlamaIndex Workflows & 2 & 2024 & llamaindex.ai & & & & May 26 \\
S5 & Memento-Skills & 4 & 2026 & arXiv & & & & May 26 \\
S5 & MetaGPT & 2 & 2023 & arXiv:2308.00352 & & & & May 26 \\
S5 & Microsoft Semantic Kernel & 2 & 2024 & learn.microsoft.com & & & &
May 26 \\
S5 & Multi-Agent Evolve (MAE) & 4 & 2025 & arXiv & & & & May 26 \\
S5 & OpenAI AgentKit & 4 & 2025 & openai.com & & & & May 26 \\
S5 & OpenAI Swarm & 1 & 2024 & openai.com & & & & May 26 \\
S5 & OpenHands / CodeAct & 3 & 2024 & arXiv:2407.16741 & & & & May 26 \\
S5 & OptiMAS & 4 & 2026 & arXiv:2608.21918 & Auto & None & C2 & Oct
26 \\
S5 & PluginRSI & 4 & 2026 & arXiv:2609.32423 & Auto & None & C2 & Oct
26 \\
S5 & Pydantic AI & 2 & 2024 & pydantic.dev & & & & May 26 \\
S5 & Recursive Harness Self-Improvement (RHI) & 4 & 2026 &
arXiv:2607.15524 & Auto & None & C3 & Oct 26 \\
S5 & Skill-MAS (meta-skill) & 4 & 2026 & arXiv:2606.18837 & Auto & None
& C2 & Oct 26 \\
S5 & TACO & 4 & 2026 & arXiv:2604.19572 & Auto & None & C2 & Oct 26 \\
S5 & Task-CoEvolve & 4 & 2026 & arXiv:2608.20169 & Auto & None & C2 &
Oct 26 \\
S5 & TOCOMAS & 4 & 2026 & arXiv:2609.37953 & Auto & None & C2 & Oct
26 \\
S5 & Vertex AI Agent Builder & 4 & 2024 & cloud.google.com & & & & May
26 \\
S5 & Workflow-Designing Agents (WDA) & 4 & 2026 & arXiv:2609.32458 &
Auto & None & C2 & Oct 26 \\
S6 & Agenta & 3 & 2024 & agenta.ai & & & & May 26 \\
S6 & Agentic Configuration Management ACM (Quessada-V & 3 & 2026 &
arXiv:2608.11166 & None & Human & n/a & Oct 26 \\
S6 & AGENTS.md / AAIF & 5 & 2025 & agents.md & & & & May 26 \\
S6 & Anthropic skills repo & 4 & 2025 & github.com/anthropics/skills & &
& & May 26 \\
S6 & Azure API Center MCP registry & 3 & 2026 & learn.microsoft.com
(ms.date & Human & Human & - & Oct 26 \\
S6 & Braintrust & 5 & 2024 & braintrust.dev & & & & May 26 \\
S6 & Cleric.ai & 2 & 2025 & cleric.ai & & & & May 26 \\
S6 & Competence Not Accuracy & 2 & 2026 & arXiv:2608.18719 & None & None
& - & Oct 26 \\
S6 & Constitutional AI & 2 & 2022 & Anthropic & & & & May 26 \\
S6 & Cursor Rules & 3 & 2024 & cursor.com & & & & May 26 \\
S6 & Decagon & 2 & 2026 & decagon.ai & & & & May 26 \\
S6 & DSPy compile & 3 & 2024 & dspy.ai & & & & May 26 \\
S6 & EvolveNet & 4 & 2026 & arXiv:2608.04968 & Auto & Auto & C2 & Oct
26 \\
S6 & EvoUndo & 4 & 2026 & arXiv:2608.28363 & Auto & None & C1 & Oct
26 \\
S6 & From Prompts to Contracts & 2 & 2026 & arXiv:2607.08028 & None &
None & C1 & Oct 26 \\
S6 & Future AGI prompt lifecycle & 4 & 2026 & futureagi.com/blog
(2026-03- & Auto & Auto & C3 & Oct 26 \\
S6 & GitHub Copilot enterprise & 4 & 2026 & github.com & & & & May 26 \\
S6 & Governance conversion (Davis et al.~case study) & 3 & 2026 &
arXiv:2607.01087 & Human & None & C2 & Oct 26 \\
S6 & Governed Capability Evolution (Qin et al.) & 3 & 2026 &
arXiv:2604.08059 & None & Auto & C2 & Oct 26 \\
S6 & HackProbe & 3 & 2026 & arXiv:2609.04665 & Auto & None & C1 & Oct
26 \\
S6 & Inspect AI & 3 & 2024 & ukgovernmentbeis.github.io & & & & May
26 \\
S6 & LangFuse & 3 & 2024 & langfuse.com & & & & May 26 \\
S6 & LangSmith Hub & 5 & 2024 & smith.langchain.com & & & & May 26 \\
S6 & MCP Gateway \& Registry (agentic-community) & 3 & 2026 &
github.com/agentic-community & Human & Human & C1 (s & Oct 26 \\
S6 & NanoResearch & 2 & 2026 & arXiv:2605.10813 & & & & May 26 \\
S6 & Official MCP Registry & 2 & 2026 & modelcontextprotocol.io/regi &
Human & None & - & Oct 26 \\
S6 & OpenAI AgentKit Evals & 3 & 2025 & openai.com & & & & May 26 \\
S6 & OPRO & 2 & 2023 & arXiv & & & & May 26 \\
S6 & Organizational Control Layer (Shi et al.) & 2 & 2026 &
arXiv:2606.04306 & None & None & n/a & Oct 26 \\
S6 & Registry-Governed Agent Lifecycle / EDDOps on AW & 3 & 2026 &
arXiv:2607.00345 & None & Auto+Human & C3 & Oct 26 \\
S6 & Reward-Free Pairwise Validator & 4 & 2026 & arXiv:2607.14408 & Auto
& None & C3 & Oct 26 \\
S6 & Runtime Skill Audit & 2 & 2026 & arXiv:2606.11671 & None & None & -
& Oct 26 \\
S6 & Runtime-Independent Persistent Agents (Enoch) & 3 & 2026 &
arXiv:2609.00546 & Human & Human & C1 & Oct 26 \\
S6 & SafeEvolve & 4 & 2026 & arXiv:2608.12851 & Auto & None & C3,C2 &
Oct 26 \\
S6 & SAGE statistical acceptance gate (≠ SAGE 2512.17 & 4 & 2026 &
arXiv:2609.36043 & Auto & None & C2 & Oct 26 \\
S6 & Self-Refine & 1 & 2023 & Madaan et al. & & & & May 26 \\
S6 & Sierra Ghostwriter & 2 & 2026 & sierra.ai & & & & May 26 \\
S6 & SkillCoach & 3 & 2026 & arXiv:2607.01874 & Auto & None & C3,C1 &
Oct 26 \\
S6 & SkillProbe & 2 & 2026 & arXiv:2603.21019 & None & None & - & Oct
26 \\
S6 & SkillSpec (intent-masked correctness) & 2 & 2026 & arXiv:2609.06052
& None & None & - & Oct 26 \\
S6 & TextGrad & 2 & 2024 & arXiv & & & & May 26 \\
S6 & TrueFoundry MCP Gateway & 3 & 2026 & truefoundry.com/docs (mcp-ga &
Human & Human & - & Oct 26 \\
S6 & Vellum & 5 & 2024 & vellum.ai & & & & May 26 \\
S6 & Verifier-as-Gatekeeper (VaG) & 4 & 2026 & arXiv:2608.05810 & Auto &
None & C2 & Oct 26 \\
S6 & W\&B Weave (Prompts) & 3 & 2024 & weave-docs.wandb.ai & & & & May
26 \\
\end{longtable}
}

\endgroup

\section{S2. Data appendix}\label{s2.-data-appendix}

\textbf{Corpus size.} The survey scores 130 unique LLM systems across
the six substrates (140 substrate-level rows in
\texttt{supplementary/p3\_master\_scores.csv}). Each system may be
scored on multiple substrates; the per-substrate row count therefore
exceeds the unique-system count.

\textbf{Core / contrast split.} 99 of the 130 unique systems satisfy the
patch-local discriminator (score ≥ 3 on at least one substrate) and form
the core corpus. The remaining 31 systems form the contrast corpus:
commonly described as agentic or self-improving in public discourse but
failing at least one discriminator. Appendix D enumerates representative
contrast systems.

\textbf{Inclusion criteria.} Any system that (a) wraps a foundation
model with at least one persistent artifact updated from deployment
signal, or (b) is named in the public discourse as a ``self-improving,''
``scaffolded,'' or ``agentic'' system. The contrast set (Self-Refine,
ReAct, TextGrad, OPRO, engineer-authored convention files; Appendix D.1)
is included so the discriminator has something to mark against; Appendix
D.2 separately lists core-corpus systems that fail the composite
two-loop criterion.

\textbf{Exclusion criteria.} Systems whose only update mechanism is
weight-level fine-tuning or RLHF: these are model-side, not
scaffold-side, and belong to the self-evolution literature this paper
differentiates from. Systems with no public evidence of deployment
(white papers without code, demos, or production claim).

\textbf{Evidence tiers.} Each system carries one or more of: (T1)
peer-reviewed publication; (T2) arXiv preprint; (T3) production claim
with metrics; (T4) open-source artifact; (T5) industry-blog or
interview-level documentation. T1/T4 carry highest weight in cluster
analysis; T3 is admissible where the claim is concrete (specific
benchmark, named integration target); T5 is admissible only when
corroborated by T1--T4 elsewhere. Tier flags are propagated to cluster
claims through the per-row evidence column in Appendix A; T3 production
claims are not pooled with T1/T2 evidence when computing cluster-level
statistics.

\textbf{Evidence-tier distribution.} Of the core corpus, approximately
one quarter carry T1 evidence, roughly half carry T2 evidence, and the
remainder carry T3 or T4 evidence as the primary tier.

\textbf{Survey window and dating.} January 2023 -- May 2026 for
systematic corpus assembly and scoring, extended by a targeted
literature refresh through September 2026 that added the systems marked
§ in Appendix A (FederatedSkill, SkillClaw, ChronoMem, SkillsVote,
HARNESSEVO) and the supporting references of §§8--10. The refresh
searched for work that would fill or narrow the §9 corner, not for
corpus completeness, so the per-substrate statistics of §6 are computed
on the May-2026 corpus. Preprint-only systems carry T2 evidence and are
flagged in Appendix A. A fraction of T2 work is recent enough that
ablations may yet be revised; cluster-level claims are robust to
dropping any single T2 system.

\textbf{Scoring discipline.} Each system received a
substrate-by-substrate 0--5 score against the rubric of Section V of the
main text and an integrated patch-local-adaptation score (PP) in
\{0,1,2,3,3.5,4,4.5,5\}. Half-scores were used sparingly when one
criterion was clearly met and another partially. Ambiguous cases were
resolved conservatively (downward). Scoring is single-rater; Appendix A
and Appendix D together expose representative per-system scores and
exclusion reasons so readers can audit ambiguous classifications
independently.

\textbf{Composite two-loop criterion (sensitivity analysis).} The strict
count of fourteen research systems satisfying the composite-two-loop
criterion changes under two natural relaxations. Relaxation (a) counts
any of versioning, lineage, or rollback as Loop-2 governance; the count
rises to approximately eighteen by re-admitting MAE, MetaGen,
MetaReflection, and CLIN. Relaxation (b) admits session signal that
updates optimizer state without persisting an artifact distinct from
optimizer state; the count rises to approximately sixteen by
re-admitting DSPy and TextGrad. The qualitative findings (Auto×Auto
well-populated; {[}Human × Human{]} and {[}Human × Auto{]} empty in the
surveyed corpus; governance-first cluster in {[}None × Auto{]}; the
cross-tenant governed-optimization cell empty) survive both relaxations.

\begin{center}\rule{0.5\linewidth}{0.5pt}\end{center}

\section{S3. The mechanism stack and continual-learning
notes}\label{s3.-the-mechanism-stack-and-continual-learning-notes}

\subsection{S3.1 The mechanism stack}\label{s3.1-the-mechanism-stack}

The abstract update rule of Section III expands into six operations that
every patch-local system must perform, whether explicitly engineered or
accidentally emergent.

\begin{enumerate}
\def\labelenumi{\arabic{enumi}.}
\tightlist
\item
  \textbf{Residual capture.} Log failures, corrections, rejected
  outputs, user edits, eval regressions, tool-call failures. The Loop-1
  signal source. Without it, \(z_t\) has no provenance.
\item
  \textbf{Mode discovery and credit assignment.} Cluster repeated
  failures into patch-specific error modes (the \emph{Architecture of
  Errors} catalog \cite{arbuzov2026architecture}, instantiated locally
  per patch) and attribute each mode to the scaffold coordinate
  responsible for it (E.6).
\item
  \textbf{Delta synthesis.} Convert a recurring mode into a candidate
  scaffold edit: an instruction rewrite, a new skill, a memory
  correction, a tool binding, an orchestration change, an eval case.
  This produces \(z_t\).
\item
  \textbf{Delta validation.} Estimate whether the edit reduces \(L_D\)
  without collateral damage on adjacent inputs. The gate \(G_t\) is
  Auto-gated when the check is an algorithmic criterion, Human-gated
  when it is reviewer judgment.
\item
  \textbf{Promotion control.} Decide whether the accepted edit stays in
  the local fork or climbs the scope hierarchy: project, stack,
  organization, cross-tenant. Loop 2.
\item
  \textbf{Drift and bloat control.} Expire stale rules, prune
  conflicting memories, roll back regressed skills, detect over-fitting.
  Without this, accumulated edits produce the scaffold-level analog of
  plasticity loss.
\end{enumerate}

Each operation is the locus of a separate engineering discipline.
Surveyed systems implement them partially and unevenly: research
full-auto systems excel at (3) and (4); production governance leaders at
(5); industry hybrids at most of (1)--(5) but rarely (6). Implementation
patterns are deferred to follow-on practitioner work.

\subsection{S3.2 Loop 1 and Loop 2 as descent steps of different
radius}\label{s3.2-loop-1-and-loop-2-as-descent-steps-of-different-radius}

Loop 1 proposes and validates local descent steps on a single \(L_D\).
Loop 2 expands the update radius. It selects, merges, and promotes
accepted local edits so that they apply to a larger averaging set of
patch losses. USER-level acceptance descends on a single \(L_D\).
PROJECT promotion descends on a project-averaged loss. STACK promotion
widens the average further. CORE or cross-tenant promotion descends on a
federated average. The two-loop design space is therefore not just a
taxonomy of engineering patterns: it is a taxonomy of \emph{which loss
each surveyed system is descending on}, gated by which estimator.

\subsection{S3.3 Overfitting in scaffold
adaptation}\label{s3.3-overfitting-in-scaffold-adaptation}

In weight training, overfitting is usually a defect. In scaffold
adaptation, local overfitting is partly the point. Frontier weights are
trained to preserve cross-patch generality; they must smooth over local
conventions and rare workflow-specific failures. A deployment patch
needs the opposite: selective specialization to local residuals. The
scaffold is therefore an intentionally overfittable layer, but one whose
overfitting is scoped, inspectable, versioned, gated, and reversible. A
scaffold that improves one repository, hospital workflow, or legal
diligence process may degrade performance elsewhere. That is not a
contradiction; it is patch adaptation. The mistake is promoting that
edit beyond its evidence radius.

In non-stationary patches (\(D = D_t\)), the same descent mechanism
tracks a moving optimum rather than converging once. This is why recency
weighting, replay, pruning, expiry, and rollback are not optional
hygiene; they are the scaffold-level analogs of continual-learning
machinery.

\begin{center}\rule{0.5\linewidth}{0.5pt}\end{center}

\section{S4. Contrast set, matrix membership, and closest approaches
(May 2026
corpus)}\label{s4.-contrast-set-matrix-membership-and-closest-approaches-may-2026-corpus}

The October 2026 additions are placed in the cells of Section VII of the
main text and listed in S1; the per-system discussion below is the May
2026 audit on which the main text's counts build.

\subsection{S4.1 Contrast set}\label{s4.1-contrast-set}

Appendix D serves two audits. \textbf{D.1} lists representative members
of the contrast corpus (31 systems that fail the patch-local
discriminator, score \textless{} 3 on every substrate), so the
discriminator has something to mark against. \textbf{D.2} lists
core-corpus systems (score ≥ 3 on at least one substrate) that
nonetheless fail the stricter composite two-loop criterion of Section V,
so the asterisked entries of the Section VII matrix are auditable.
Systems named in the text but not scored as rows in
\texttt{supplementary/p3\_master\_scores.csv} (ReAct, MetaGen,
MetaReflection, CLIN, ACE) are classified from their papers and marked
with a dagger.

\subsubsection{S4.1.1 Contrast corpus: fails the patch-local
discriminator}\label{s4.1.1-contrast-corpus-fails-the-patch-local-discriminator}

{\def\LTcaptype{none} % do not increment counter
\begin{longtable}[]{@{}
  >{\raggedright\arraybackslash}p{(\linewidth - 6\tabcolsep) * \real{0.2500}}
  >{\raggedright\arraybackslash}p{(\linewidth - 6\tabcolsep) * \real{0.2500}}
  >{\raggedright\arraybackslash}p{(\linewidth - 6\tabcolsep) * \real{0.2500}}
  >{\raggedright\arraybackslash}p{(\linewidth - 6\tabcolsep) * \real{0.2500}}@{}}
\toprule\noalign{}
\begin{minipage}[b]{\linewidth}\raggedright
System
\end{minipage} & \begin{minipage}[b]{\linewidth}\raggedright
Year
\end{minipage} & \begin{minipage}[b]{\linewidth}\raggedright
Discriminator failure
\end{minipage} & \begin{minipage}[b]{\linewidth}\raggedright
Brief reason
\end{minipage} \\
\midrule\noalign{}
\endhead
\bottomrule\noalign{}
\endlastfoot
Self-Refine \cite{madaan2023selfrefine} & 2023 & No persistent artifact
& In-context iteration only \\
ReAct \cite{yao2023react}† & 2023 & No persistent artifact &
Thought-act-observe loop is intra-session \\
TextGrad \cite{yuksekgonul2024textgrad} & 2024 & No artifact distinct
from optimizer state & Textual-gradient optimization; SkillOpt
{[}\cite{yang2026skillopt}; Section III.5; Appendix G{]} adds the
persistent \texttt{best\_skill.md}, held-out validation gate,
rejected-edit memory, and textual learning-rate schedule that TextGrad
lacks. \\
OPRO, Toolformer & 2023 & No persistent deployment artifact &
Optimization or tool use inside a single training or inference run \\
Aider conventions, Lovable knowledge & 2024 & No feedback loop &
Engineer-authored convention files \\
AGENTS.md (as a standard) & 2025 & No feedback loop & Cross-vendor
artifact format; adoption measured, adaptation absent (Section VI) \\
LangChain Skills, HuggingFace Agent Skills, CrewAI Skills & 2025--2026 &
No feedback loop & Curated skill packs without failure-triggered
update \\
Tran \& Kiela \cite{tran2026singleagent} & 2026 & Not a system &
Single-agent baselines match or beat multi-agent under matched
thinking-token budgets; constrains S5 claims (see Section VI) \\
\end{longtable}
}

\subsubsection{S4.1.2 Core corpus, fails the composite two-loop
criterion}\label{s4.1.2-core-corpus-fails-the-composite-two-loop-criterion}

{\def\LTcaptype{none} % do not increment counter
\begin{longtable}[]{@{}
  >{\raggedright\arraybackslash}p{(\linewidth - 6\tabcolsep) * \real{0.2500}}
  >{\raggedright\arraybackslash}p{(\linewidth - 6\tabcolsep) * \real{0.2500}}
  >{\raggedright\arraybackslash}p{(\linewidth - 6\tabcolsep) * \real{0.2500}}
  >{\raggedright\arraybackslash}p{(\linewidth - 6\tabcolsep) * \real{0.2500}}@{}}
\toprule\noalign{}
\begin{minipage}[b]{\linewidth}\raggedright
System
\end{minipage} & \begin{minipage}[b]{\linewidth}\raggedright
Year
\end{minipage} & \begin{minipage}[b]{\linewidth}\raggedright
Composite-criterion failure
\end{minipage} & \begin{minipage}[b]{\linewidth}\raggedright
Brief reason
\end{minipage} \\
\midrule\noalign{}
\endhead
\bottomrule\noalign{}
\endlastfoot
Reflexion \cite{shinn2023reflexion} & 2023 & No Loop 2 & Per-agent
episodic memory without promotion \\
Generative Agents & 2023 & No Loop 2 & Per-agent reflection + importance
memory \\
ExpeL & 2023 & Loop boundary blurred & Experience bank built from
training tasks \\
Mem0 (vanilla) & 2024 & No Loop 2 & Per-user memory; no cross-context
promotion \\
Zep & 2024 & No Loop 2 & Per-agent/session knowledge graph \\
ChatGPT Memory & 2024 & No Loop 2 & Per-user persistence \\
Claude Memory & 2025 & No Loop 2 & Per-conversation-thread
persistence \\
MemGPT & 2023 & No Loop 2 & Virtual-context management per agent \\
DSPy & 2023 & No artifact distinct from optimizer state &
Prompt-template optimization; SkillOpt {[}\cite{yang2026skillopt};
Section III.5; Appendix G{]} realizes the artifact-bearing form. \\
MetaGen† & 2026 & No Loop 2; no cross-session persistence & Role pool
dynamic within single inference call \\
MAE (Multi-Agent Evolve) & 2025 & No Loop 2 & Co-evolution within
Proposer/Solver/Judge triplet \\
MetaReflection† & 2024 & No Loop 2 & Per-agent semantic memory \\
CLIN† & 2023 & No cross-instance promotion & Causal abstractions
per-agent-per-environment \\
LangGraph, AutoGen, CrewAI & 2023--2024 & Framework only & Orchestration
frameworks; not deployed systems \\
W\&B Weave, Agenta, Langfuse & 2024 & No Loop 1 & Observability /
eval-as-code \\
Mobile-Agent-E, Cradle, OS-Copilot & 2024--2025 & No Loop 2 &
Self-evolving but per-session/per-deployment \\
Agent S, Agent S2 & 2024--2025 & No Loop 2 & Computer-use agentic
framework \\
ACE \cite{ace2025agentic}† & 2025 & No Loop 2 in surveyed form &
Single-context evolving playbook \\
Cursor Rules (vanilla) & 2024 & No Loop 2 & Per-project artifacts \\
GitHub Copilot custom instructions, Windsurf, Continue.dev, Zed AI rules
& 2024 & No Loop 1 & Engineer-authored at repo base branch \\
Replit \texttt{replit.md} & 2025 & Loop 2 absent & Self-writing scaffold
but per-workspace \\
\end{longtable}
}

The entries above cover the main exclusion families (in-context-only; no
feedback loop; not a system) and the main composite failures (no Loop 2;
framework-only; engineer-authored without failure trigger;
optimizer-state-only); finer-grained per-system breakdowns are deferred
to follow-on practitioner work.

\begin{center}\rule{0.5\linewidth}{0.5pt}\end{center}

\subsection{S4.2 Survey detail}\label{s4.2-survey-detail}

This appendix carries the cell-by-cell discussion of the two-loop matrix
from Section VII, the closest-approaches enumeration that supports
Section IX, and the deployment-dependence argument. It supplies the
verification material for §Section VII--9.

\subsubsection{S4.2.1 Two-loop matrix: cell-by-cell
discussion}\label{s4.2.1-two-loop-matrix-cell-by-cell-discussion}

The named membership of the Section VII matrix (Figure 3):

{\def\LTcaptype{none} % do not increment counter
\begin{longtable}[]{@{}
  >{\raggedright\arraybackslash}p{(\linewidth - 6\tabcolsep) * \real{0.2500}}
  >{\raggedright\arraybackslash}p{(\linewidth - 6\tabcolsep) * \real{0.2500}}
  >{\raggedright\arraybackslash}p{(\linewidth - 6\tabcolsep) * \real{0.2500}}
  >{\raggedright\arraybackslash}p{(\linewidth - 6\tabcolsep) * \real{0.2500}}@{}}
\toprule\noalign{}
\begin{minipage}[b]{\linewidth}\raggedright
\end{minipage} & \begin{minipage}[b]{\linewidth}\raggedright
\textbf{Loop-2 Auto-gated}
\end{minipage} & \begin{minipage}[b]{\linewidth}\raggedright
\textbf{Loop-2 Human-gated}
\end{minipage} & \begin{minipage}[b]{\linewidth}\raggedright
\textbf{Loop-2 None}
\end{minipage} \\
\midrule\noalign{}
\endhead
\bottomrule\noalign{}
\endlastfoot
\textbf{Loop-1 Auto-gated} & NanoResearch, SkillOpt, SkillRL,
AlphaEvolve, ADAS, DGM, Voyager, SAGE, EvolveR, CoEvoSkills, AutoAgent,
AutoManual & SkillForge, CASCADE, Sierra Agent OS 2.0†, Cognition Devin†
& MAE*, MetaGen*, MetaReflection*, CLIN*, DSPy*, TextGrad* \\
\textbf{Loop-1 Human-gated} & \textbf{{[}empty by engineering logic{]}}
& \textbf{{[}empty in surveyed corpus{]}} & \emph{(none)} \\
\textbf{Loop-1 None} & ExpeL (weak), Braintrust*, LangSmith Hub*,
Vellum* & Anthropic skills repository* & Self-Refine*, ReAct*, Mem0* \\
\end{longtable}
}

\emph{Asterisks fail at least one composite-two-loop criterion; included
for contrast.} † Industry exemplar evaluated using vendor documentation
as primary evidence. DSPy and TextGrad sit in {[}Auto × None{]} because
their Loop 1 is automated but persists only optimizer state (Appendix B,
relaxation (b)). The matrix is restricted to the 22 systems of the
strict composite audit plus SkillOpt; Appendix A also assigns Loop-1 /
Loop-2 entries to systems outside that audit (COSPLAY, Cuadros et al.,
PACE, CORAL, GEPA, Meta-Harness, MOSS, RCL, marked ‡ there, and the
September-2026 refresh additions marked §), but those entries are read
from the papers and were not re-audited under the composite criterion,
so they are neither placed in the matrix nor counted in Section V.

\textbf{{[}Auto × Auto{]}} is most populated (twelve systems). Every
member has a machine-evaluated decision criterion on both loops.
Velocity is the cluster's defining property; the lack of
cross-organizational governance is its uniform weakness.

\textbf{{[}Auto × Human{]}} is the production-viable cell: Loop 1 fires
at machine speed under an algorithmic criterion, Loop 2 requires human
approval before propagation. The pattern reconciles automation velocity
with reviewer accountability. SkillForge's LLM-judge + reviewer pipeline
is the architectural exemplar. CASCADE, Sierra Agent OS 2.0, and Devin
instantiate variants; the industry exemplars (Agent OS 2.0, Devin) are
marked with † because vendor documentation rather than peer-reviewed
metrics is the primary evidence.

\textbf{{[}None × Auto-gated{]}} holds the governance-first
prompts-as-code systems (Braintrust, LangSmith Hub, Vellum). They have a
real Loop 2 (eval-gated CI blocks promotion on regression), but no Loop
1 in the patch-local sense, because the candidate artifacts are authored
by engineers at their desks rather than triggered by session feedback.
The asterisks mark this: governance without inner-loop adaptation. The
cell is informative because it shows what survives when only the outer
loop is automated.

\textbf{{[}Human × Auto{]}} is empty by engineering logic: if humans
author candidates manually, automated promotion offers little leverage.
A speculative occupant would be ``expert curation + auto-distribution''
(expert-authored skills automatically gated through a held-out eval
before shipping); no published example exists.

\textbf{{[}Human × Human{]}} is empty in the surveyed corpus. A
candidate occupant would be a reviewer-authored-and-reviewer-promoted
artifact pipeline targeting solo and small-team workflows where
automation cost exceeds review latency. No public example was found in
the survey window.

\textbf{{[}Human × None{]}} is structurally improbable: a system
requiring human authorship on Loop 1 but with no Loop 2 collapses to
``engineer-edits-a-config'' and is not patch-local in the sense of
Section III.

\textbf{DGM dual-mode note.} DGM operates in both {[}Auto × Auto{]}
(primary archive-based evolution) and {[}Auto × Human{]} (sandboxed
code-modification review) modes depending on deployment configuration;
we list it once in the dominant cell.

\subsubsection{S4.2.2 Closest approaches to governed cross-tenant
scaffold
optimization}\label{s4.2.2-closest-approaches-to-governed-cross-tenant-scaffold-optimization}

The four properties of governed cross-tenant scaffold optimization are
(a) auto-gated Loop 1, (b) multi-tenant cross-org promotion, (c)
versioned lineage with RBAC, (d) rollback. Closest approaches partition
the requirements:

\begin{itemize}
\tightlist
\item
  \textbf{AlphaEvolve} has (a) and partial (b) within a single
  organization; no cross-org promotion mechanism.
\item
  \textbf{NanoResearch} has (a) at the research level; no enterprise
  governance layer described.
\item
  \textbf{SkillOpt} has (a) fully and partial (c) via the exported
  \texttt{best\_skill.md} artifact plus protected slow-update field;
  reports positive cross-model, cross-harness, and cross-benchmark
  transfer as empirical evidence-radius measurement; does not implement
  (b) cross-organizational promotion or (d) explicit rollback. The
  Appendix G walkthrough maps every SkillOpt design decision to its
  Section III counterpart.
\item
  \textbf{CORAL} \cite{qu2026coral} has partial (a) via heartbeat-gated
  promotion and is the only system in the corpus that implements
  asynchronous multi-agent Loop 2 within one organization through shared
  persistent memory; no (b) cross-organizational scope, no formal (c)
  RBAC, no (d) rollback.
\item
  \textbf{PACE} \cite{ling2026pace} has (a) at the small-model
  production timescale and a fast/slow two-timescale Loop 1 + Loop 2
  that consolidates across one deployment's own episodes rather than
  across deployments; no (b), no formal (c) or (d). PACE is the closest
  demonstration that the two-loop architecture is implementable under
  small-frozen-model constraints.
\item
  \textbf{Meta-Harness} \cite{lee2026metaharness} has (a) over the
  harness substrate (S5) with a filesystem-of-candidates
  credit-assignment record that functions as the accumulated history
  \(H_t\) of Section III; no Loop-2 cross-organizational promotion, no
  RBAC, no rollback.
\item
  \textbf{SkillForge} has (a) and (b) within one enterprise, partial (c)
  via VFS commits; (d) is undocumented in the public report.
\item
  \textbf{Sierra Agent OS 2.0} has (b) and partial (a); the GitHub-style
  Workspaces provide some lineage but not RBAC-with-rollback in the
  strict sense.
\item
  \textbf{FederatedSkill} \cite{yang2026federatedskill} has (a) and the
  strongest partial (b) in the corpus: heterogeneous clients
  (backbone-plus-harness profiles under one federation operator)
  exchange semantic skill diffs, and a server-side evolution agent
  applies, adapts, or rejects each patch per client; partial (c) via an
  auditable decision log recording source patch, reward, and
  justification per modified path; no RBAC, no (d).
\item
  \textbf{SkillClaw} \cite{ma2026skillclaw} has partial (a) via a
  nightly validation step and partial (b) across users of one shared,
  synchronized skill repository; no (c) or (d) reported.
\item
  \textbf{ChronoMem} \cite{su2026chronomem} has (d) and the versioning
  half of (c) for the memory substrate alone (whole-memory snapshots,
  version histories, natural-language rollback inside Google's Agent
  Development Kit); no Loop 1, no (b), no RBAC.
\item
  \textbf{Vellum / Braintrust / LangSmith Hub} have (c) and (d)
  explicitly but no (a) or (b) in the strict sense; these are
  governance-first systems by construction.
\end{itemize}

No surveyed system combines (b), (c), and (d) under any gate type. The
cross-organizational governance triad is itself unfilled, independently
of whether Loop 1 is auto-gated.

\subsubsection{S4.2.3 Deployment-dependence of the missing
architecture}\label{s4.2.3-deployment-dependence-of-the-missing-architecture}

Whether the corner \emph{should} be filled is deployment-dependent. In
safety-critical domains (medical decision support, legal advice,
financial advice), reviewer-judgment gates may be a design feature
rather than a defect, and the human latency is the point. The narrower
survey claim is what matters: no public system currently combines
automated local scaffold evolution with governed cross-organization
promotion, versioned lineage with RBAC, and rollback. Naming the corner
converts a vague gap into a specific four-property audit criterion that
any future system can be measured against. The shortest demonstrable
path appears to build on AlphaEvolve's cascade evaluator and MAP-Elites
archive, extending federation from within-DeepMind multi-target to
cross-customer multi-tenant. The architectural surface, not the
timeline, is what this paper claims.

\begin{center}\rule{0.5\linewidth}{0.5pt}\end{center}

\section{S5. The convergent research families, system by
system}\label{s5.-the-convergent-research-families-system-by-system}

This section carries the May--September 2026 walkthroughs; the October
2026 additions to each family are described in Section VIII of the main
text and scored in S1.

This appendix expands Section VIII along two axes. G.1 indexes the
landscape; G.2--G.3 walk SkillOpt mechanism-by-mechanism; G.4--G.8 walk
each convergent research family, stating its Section III mapping
explicitly and naming what it leaves open; G.9 synthesizes four claims
supported by the families as a whole; G.10--G.11 carry the Section
IX-property audit of SkillOpt and its published cost-per-point figure.

\subsection{S5.1 The convergent
landscape}\label{s5.1-the-convergent-landscape}

SkillOpt \cite{yang2026skillopt} is the cleanest current instance of
artifact-layer descent in our survey, but it sits alongside at least
eighteen independent research threads in 2025--2026 that converge on the
same disciplined gradient-like update mechanism for the scaffold layer.
The convergence is striking because it happens \emph{from different
starting points}: GEPA and Tiwari et al.~start from prompt-optimization
and continual-learning respectively; RCL starts from cognitive-agent
learning; Meta-Harness, AHE, MOSS, CANTANTE, and Ong et al.~start from
harness engineering; MOCHA and Ratchet start from skill-deployment
constraints; ICT, Combee, CORAL, and PACE start from cross-task /
multi-agent / production-deployment angles. Each system instantiates a
different subset of Section III's five conditions, and each exposes a
different gap.

The landscape table of Section VIII indexes the families by Section III
mechanism and by what each leaves open; the subsections below walk each
family in detail. By substrate (Section V), GEPA and RCL operate on the
prompt and context layer (S1--S2); the harness family on orchestration
(S5), with MOSS extending to source code (S2+S5); MOCHA and Ratchet on
skills under deployment constraints (S2); the transfer family on
integrated memory-plus-orchestration stacks (S3+S5); and SkillOpt on an
S2-led integrated stack.

Read against Section IX, no member of any family combines (a) auto-gated
Loop 1 with (b) cross-organizational Loop 2, (c) versioned lineage with
RBAC, and (d) rollback. The Section III picture is realized in fragments
across the families; the Section IX gap is confirmed by every member.

\subsection{S5.2 SkillOpt: full
correspondence}\label{s5.2-skillopt-full-correspondence}

The table below extends the Section VIII compact correspondence with the
Section III role of each component and an algorithm-level pointer into
SkillOpt's published procedure (their Algorithm 1).

{\def\LTcaptype{none} % do not increment counter
\begin{longtable}[]{@{}
  >{\raggedright\arraybackslash}p{(\linewidth - 6\tabcolsep) * \real{0.2500}}
  >{\raggedright\arraybackslash}p{(\linewidth - 6\tabcolsep) * \real{0.2500}}
  >{\raggedright\arraybackslash}p{(\linewidth - 6\tabcolsep) * \real{0.2500}}
  >{\raggedright\arraybackslash}p{(\linewidth - 6\tabcolsep) * \real{0.2500}}@{}}
\toprule\noalign{}
\begin{minipage}[b]{\linewidth}\raggedright
Section III quantity
\end{minipage} & \begin{minipage}[b]{\linewidth}\raggedright
SkillOpt instantiation
\end{minipage} & \begin{minipage}[b]{\linewidth}\raggedright
Section III role explained
\end{minipage} & \begin{minipage}[b]{\linewidth}\raggedright
Algorithm pointer
\end{minipage} \\
\midrule\noalign{}
\endhead
\bottomrule\noalign{}
\endlastfoot
\(\theta_s\) (scaffold) & \texttt{best\_skill.md} --- a 379--1,995-token
markdown document & the trainable external state of the frozen target
model; persistent and exported as the deployment artifact &
\(s_{\text{best}}\) (Algorithm 1, lines 1, 22, 37) \\
\(L_D\) (patch loss) & benchmark-native hard score / exact-match on a
held-out test split & scalar quality signal that the optimizer is
implicitly minimizing & \(\textsc{Evaluate}(M, h, s, D)\) (lines 2, 17,
37) \\
\(Q_t\) (proposal) & optimizer-model reflection over failure and success
minibatches; failures and successes analyzed separately, then merged
with priority on failure fixes & proposal distribution conditioned on
accumulated rollout evidence and current scaffold & optimizer calls
\(O\) for failure analysis, success analysis, and merge (lines 9--11) \\
\(z_t\) (edit) & structured add/delete/replace patch on the skill (or
rewrite-from-suggestions in rewrite mode) & the candidate update applied
via \(\oplus\) to the current scaffold & rank/clip to \(\eta_t\) edits,
apply to obtain \(\tilde s\) (lines 12--13) \\
\(\eta_t\) (textual learning rate) & bounded per-step edit budget;
constant / linear / cosine / autonomous schedules; default cosine with
floor \(\eta_t=2\) & edit-magnitude controller; the Section III
\emph{learning rate} in textual form & rank/clip step (line 12);
schedule applied per epoch \\
\(G_t\) (gate) & held-out selection-split accept gate, strictly greater
than current selection score & the per-edit gate \(G_t\) --- accepts
only when \(\widehat{\Delta}_{z_t} L_D\) improves the selection score &
comparison at line 20; cache lookup at lines 14--18 \\
rejected-edit buffer & epoch-local store of failed proposals and the
score drop they caused & converts rejected steps into negative feedback
for later proposals within the same epoch; zero deployment cost &
\(\mathcal B\) updated at line 26; reset at line 5 \\
slow/meta update & epoch-end longitudinal guidance written to a
markup-fenced protected slow-update field, still passed through
\(D_{\text{sel}}\) & second-order consolidation across adjacent epochs;
the Section III ``step-level edits cannot overwrite the protected
slow-update field'' decomposition & lines 28--31; gate applied to
injected guidance \\
optimizer-side meta skill & \(m_{\text{meta}}\) --- a prompt-side
summary of which edits helped, which failed, and which failures
persisted & training-time guidance for \(Q_t\); never shipped with the
deployed artifact & lines 33--34 (training-side only) \\
\end{longtable}
}

The mapping is one-to-one at the level of Algorithm 1: every Section III
quantity has a named, code-level counterpart, and every SkillOpt control
surface maps onto a Section III quantity. SkillOpt is therefore the
cleanest published instance we are aware of in which the Section III
picture is \emph{mechanism} rather than \emph{metaphor}: all five
Section III conditions are instantiated.

\subsection{S5.3 What SkillOpt does that prior textual-optimization
systems do
not}\label{s5.3-what-skillopt-does-that-prior-textual-optimization-systems-do-not}

SkillOpt's own experiment design compares it against six baselines under
an identical protocol (same target model, same held-out test split, same
scorer). The differentiation below frames each baseline in Section III
vocabulary and identifies the Section III mechanism the baseline lacks.
The through-line is the paper's thesis: \emph{the scaffold/evolution
layer needs a disciplined gradient-like update mechanism}, and the live
distinction across baselines is which Section III conditions each one
fails.

\begin{itemize}
\item
  \textbf{TextGrad} \cite{yuksekgonul2024textgrad}\textbf{.}
  Backpropagates textual feedback through \texttt{tg.Variable} objects
  in a compound system. In Section III terms: \(z_t\)-style edits exist,
  but \(\theta_s\) does not --- TextGrad has no artifact distinct from
  optimizer state. There is no \(\eta_t\) schedule, no held-out \(G_t\),
  and no rejected-edit buffer. SkillOpt's measured advantage over
  TextGrad on direct chat ranges from \(+5.9\) to \(+34.0\) points
  across benchmarks at GPT-5.5 scale and is uniformly positive at all
  seven target-model scales. This matches our Appendix D classification
  of TextGrad as failing the patch-local discriminator on persistence;
  SkillOpt closes precisely that gap.
\item
  \textbf{GEPA} \cite{agrawal2025gepa}\textbf{.} Reflective prompt
  evolution with Pareto-style multi-candidate selection. Closer than
  TextGrad in that it carries multiple candidates and a selection step.
  In Section III terms: \(Q_t\) is real, but the optimization target is
  a prompt rather than a persistent skill artifact, \(G_t\) is not the
  strictly-greater held-out gate of Section III, and there is no
  \(\eta_t\) bound --- accepted edits can wholesale rewrite. SkillOpt
  matches or outperforms GEPA in every reported cell of its main results
  table \cite{yang2026skillopt}; we do not restate the per-cell deltas
  here.
\item
  \textbf{Trace2Skill} \cite{ni2026trace2skill}\textbf{.} Distils
  trajectory-local lessons into reusable skill artifacts. In Section III
  terms: \(\theta_s\) exists, but there is no iterative \(G_t\) ---
  Trace2Skill mines skills offline rather than running a held-out
  validation loop that accepts or rejects each candidate. Our Section
  III mechanism stack therefore loses \emph{delta validation} (operation
  4) and the \emph{rejected-edit buffer}. SkillOpt's measured gap over
  Trace2Skill is \(+5.0\) to \(+19.0\) points across benchmarks at
  GPT-5.5 direct chat.
\item
  \textbf{EvoSkill} \cite{alzubi2026evoskill}\textbf{.} Skill-folder
  evolution under failure analysis, deployed harness-side. The strongest
  baseline on the Codex and Claude Code harness rows. In Section III
  terms: \(\theta_s\) is a folder rather than a single artifact, and
  credit assignment via failure analysis is genuine, but there is no
  \(\eta_t\) bound on per-step folder edits and no rejected-edit buffer.
  SkillOpt beats EvoSkill by \(+14.0\) inside Codex and \(+3.2\) inside
  Claude Code on the GPT-5.5 five-benchmark average.
\item
  \textbf{Human skill} (expert-written, 145--516 tokens per benchmark).
  In Section III terms: \(\theta_s\) exists, but Loop 1 does not --- the
  artifact is static at deployment. Provides the no-Loop-1 ceiling.
\item
  \textbf{One-shot LLM skill} (generated zero-shot from a benchmark
  description by GPT-5.5, never updated). In Section III terms:
  \(\theta_s\) exists; Loop 1 does not; \(Q_t\) fires exactly once with
  no rollout evidence. The ``untrained scaffold'' baseline.
\end{itemize}

ABSTRAL \cite{song2026abstral} and EvoTest \cite{he2025evotest},
referenced in SkillOpt's related work but not measured in their table,
operate at the multi-agent design / test-time configuration layer rather
than the persistent skill layer; they are out of scope for the
artifact-layer-descent comparison.

Synthesizing across all six measured baselines: the live difference is
whether the loop has (i) a persistent edit-target distinct from
optimizer state (\(\theta_s\)), (ii) a held-out validation gate on each
accept (\(G_t\) with \(\widehat\Delta\) estimator and accept threshold
\(\tau_t\)), (iii) negative-feedback retention (rejected-edit buffer
feeding \(Q_t\)), and (iv) a bounded learning rate with a schedule
(\(\eta_t\) + cosine/linear/constant/autonomous). SkillOpt is the first
published system in our survey that instantiates all four
simultaneously, and it does so as a single offline training loop whose
deployed artifact remains a 379--1,995-token markdown file requiring no
inference-time optimizer calls.

A follow-up, SkillOpt-Lite \cite{shen2026skilloptlite}, asks which of
these controls are necessary. It reformulates skill optimization as
zeroth-order optimization, keeps three principles (file-system-based
trajectory exploration, consensus attribute mining that extracts
attributes common across heterogeneous rollouts and discards
single-sample eccentricities, and an independent held-out validation
gate that is strictly disjoint from the training failures), and removes
minibatch reflection pooling, the rejected-edit buffer, the slow/meta
update, and the textual learning-rate schedule. It reports faster
convergence and higher scores than full SkillOpt (LiveMath +8.8 points
at GPT-5.5 and +25.4 at GPT-5.4-nano). Read against the four conditions
above, this is evidence that (ii) the gate and a consensus form of
credit assignment are load-bearing, while (iii) negative-feedback
retention and (iv) the bounded learning rate are stabilizers whose value
is benchmark-dependent; the finding does not touch the Section IX
governance properties.

\subsection{S5.4 GEPA, Tiwari, and the reflective-evolution
family}\label{s5.4-gepa-tiwari-and-the-reflective-evolution-family}

The immediate predecessor to SkillOpt's optimizer framing is
\textbf{GEPA} \cite{agrawal2025gepa}. It shows that natural-language
reflection over execution trajectories can provide a richer adaptation
signal than sparse scalar rewards for prompt evolution in compound AI
systems. GEPA samples trajectories containing reasoning steps, tool
calls, and tool outputs; reflects on them to diagnose failures and
propose prompt updates; and maintains a Pareto frontier of candidate
prompts rather than greedily mutating only the current best. Across six
tasks, GEPA outperforms GRPO by 6\% on average and by up to 20\% while
using up to 35× fewer rollouts, and also outperforms MIPROv2 in the
reported aggregate comparison. In Section III vocabulary, GEPA strongly
instantiates the proposal side, \(Q_t\), through trajectory-conditioned
reflection, and uses Pareto-based candidate selection as a
proposal/pruning mechanism. What it does not instantiate is the full
SkillOpt-style artifact lifecycle: a bounded edit budget on a reusable
skill document, a strictly-greater held-out gate on each accepted skill
update, a rejected-edit buffer used as negative feedback, and an
exported deployment artifact with governance hooks.

\textbf{\cite{tiwari2026fastandslow}} supply the broader theoretical
framing from the optimization side, arguing explicitly that restricting
learning to \emph{either} in-context (prompt optimization) \emph{or}
in-weights (RL / fine-tuning) is unnecessary and costly. In-context
adaptation is cheap and rapid but capacity-bounded; in-weights
adaptation is expressive but induces catastrophic forgetting and loss of
plasticity \cite{dohare2024loss}. The scaffold layer of Section II is
precisely the substrate that absorbs what in-context adaptation can
carry cheaply, leaving the weights free to generalize. Tiwari et
al.~supply the two-timescale motivation from the optimization side;
Section II supplies the matching motivation from the
production-reliability side. The two framings are independent and
mutually reinforcing.

\subsection{S5.5 Reflective Context Learning: a parallel
formalization}\label{s5.5-reflective-context-learning-a-parallel-formalization}

The most direct independent convergence on the Section III formalization
comes from \textbf{Reflective Context Learning} {[}RCL;
\cite{vassilyev2026rcl}{]}. Their opening observation is that ``the
fundamental problems of learning, including credit assignment,
overfitting, forgetting, local optima, and high-variance learning
signals, persist whether the learned object lies in parameter space or
context space,'' and that current methods are ``fragmented and ad hoc''
precisely because they address these challenges one at a time rather
than as a unified optimization problem over the context. The parallel to
Section III's argument --- that scaffold maintenance needs the
discipline of weight-space optimization applied to an external artifact
layer --- is exact and arrived at independently.

RCL maps onto Section III as follows. The \emph{learned context} is
\(\theta_s\). The \emph{reflection signal} is a form of credit
assignment identifying which context coordinate caused a failure
(Appendix E.1, operation 2). The \emph{generalization penalty} that RCL
introduces to prevent narrow sample-specific rules is structurally the
complexity penalty \(\Omega(z)\) of the gate condition. Where RCL and
Section III differ is in Loop-2 scope: RCL is framed around single-agent
cross-task generalization, while Section III's Loop 2 is
cross-deployment, cross-tenant, and cross-organizational promotion. The
convergence is nonetheless significant: two research threads starting
from different motivations (production reliability in Section III,
cognitive agent learning in RCL) arrive at the same formal observation
that the classical learning problems do not disappear when the learned
object moves outside the weights.

\subsection{S5.6 The harness-optimization
cluster}\label{s5.6-the-harness-optimization-cluster}

A second convergent family focuses on the \emph{harness} --- the code
and configuration layer that wraps a frozen model --- rather than on
skill documents specifically. This family independently identifies
credit assignment as the central unsolved problem of scaffold evolution,
and its engineering solutions directly instantiate Appendix E.1's
operation 2 (mode discovery and credit assignment) and operation 4
(delta validation).

\textbf{Meta-Harness} \cite{lee2026metaharness} is the cleanest
representative. Its opening claim --- that ``the performance of large
language model systems depends not only on model weights, but also on
their harness: the code that determines what information to store,
retrieve, and present to the model'' --- is §1's scaffold-versus-weights
decomposition restated for the harness substrate specifically.
Meta-Harness then identifies that existing text optimizers ``compress
feedback too aggressively,'' losing the fine-grained signal needed for
credit assignment. Its solution is an outer-loop proposer with access to
source code, scores, and execution traces of all prior candidates
through a filesystem, so that the edit history itself becomes the
credit-assignment signal. In Section III terms, the filesystem of prior
candidates is a concrete form of the accumulated history \(H_t\) that
conditions the proposal distribution \(Q_t\), and the outer loop is a
Loop-1 closure over the harness substrate (Section V, S5) rather than
the skill substrate (S2).

\textbf{Agentic Harness Engineering} {[}AHE; \cite{lin2026ahe}{]}
attacks the same problem from the observability side. The core diagnosis
is that harness engineering ``faces a heterogeneous action space across
editable components, voluminous trajectories that bury actionable
signal, and edits whose effect is hard to attribute.'' These three
failure modes map precisely onto Section III's mechanism stack: a
heterogeneous action space is the coordinate-identification problem (E.1
op 2); voluminous trajectories that bury signal is the residual-capture
problem (op 1); and effects that are hard to attribute is the
credit-assignment half of op 2 (E.6). AHE's response --- three matched
observability pillars giving every editable harness component a
file-level representation --- is an engineering instantiation of the
coordinate decomposition that Section III assumes. The implicit claim is
that artifact-layer descent requires not just the optimization loop but
also the instrumentation that makes scaffold coordinates individually
observable and attributable.

\textbf{MOSS} \cite{cai2026moss} extends the harness argument by
observing that all prior self-evolving agent systems ``confine evolution
to text-mutable artifacts --- skill files, prompt configurations, memory
schemas, workflow graphs --- and leave the agent harness untouched.''
Since routing, hook ordering, state invariants, and dispatch live in
code rather than in any text artifact, an entire class of structural
failure is, in MOSS's phrase, ``physically unreachable from the text
layer.'' MOSS therefore extends Loop-1 evolution to source-level code
rewriting. In Section V terms, MOSS evolves S2 (Skills) and S5
(Orchestration) simultaneously, rather than treating the orchestration
substrate as a fixed wrapper around an evolvable skill layer. This
matters for Section IX: a governed cross-tenant scaffold optimizer that
covers only text-mutable artifacts leaves the code substrate ungoverned,
and MOSS's diagnosis suggests the gap is not merely inconvenient but
structurally incomplete.

\textbf{CANTANTE} \cite{zehle2026cantante} attacks credit assignment
directly in the multi-agent setting. Its observation is that ``scores
are available only at the system level, whereas the parameters governing
agent behavior are local,'' making optimization of multi-agent scaffolds
``fundamentally a credit-assignment problem.'' CANTANTE's solution is
contrastive: it decomposes system-level rewards into per-agent update
signals by contrasting rollouts of multiple jointly acting
configurations. In Section III vocabulary, CANTANTE provides a concrete
mechanism for the credit-assignment step (E.1 op 2; E.6) in the case
where the scaffold is distributed across agents rather than concentrated
in a single artifact. The gap it closes is the same one Section III
identifies in Appendix E.6 --- that attribution (output error assigned
to an upstream scaffold coordinate) must succeed with non-zero
probability per failure --- but instantiated for the multi-agent
topology substrate (S5).

\textbf{Harness optimizer evaluation} \cite{ong2026harnesseval}
provides independent empirical evidence for a failure mode this paper
flags in Section X: that evaluating scaffold optimizers ``solely by
observing target agents' performance gains'' misses intermediate
erroneous edits that hinder performance. Their finding --- that it is
``unclear whether harness optimization is driven by optimizers' informed
update actions or simply trial-and-error'' --- is a direct empirical
instantiation of the distinction between Stage 2 (random artifact
search) and Stage 3 (artifact-layer descent) from Section II's table.
The credit-assignment problem is not merely theoretical: production
harness optimizers that cannot distinguish informed updates from lucky
trial-and-error are operating in Stage 2 even when they report benchmark
improvements.

\textbf{HARNESSEVO} \cite{nguyen2026harnessevo} turns the same
diagnosis into a measurement. It splits the harness into four separately
evolvable slots (role, task strategy, tool and format rules, reflection
and control) and attributes value by leave-one-in and leave-one-out
evolution. On ALFWorld with a 7B target model nearly all of the gain
lives in the reflection/control slot (leave-one-in +0.119; the other
slots individually null), and dividing a fixed 64-rollout budget evenly
across the four slots leaves each below the optimizer's effective search
floor, so every slot freezes at its seed; concentrating the budget on
the high-value slot recovers the gain (0.642 to 0.761). On WebShop the
effect is absent. In Section III vocabulary this is direct evidence that
credit assignment over scaffold coordinates must precede budget
allocation, and that the coordinate decomposition of Section V is not
merely descriptive: which coordinate carries value is task-dependent and
must be measured, not assumed.

\subsection{S5.7 Multi-objective and lifecycle
constraints}\label{s5.7-multi-objective-and-lifecycle-constraints}

Two papers expose constraints on the Section III picture that
single-loss descent does not capture, both relevant to the Section IX
missing-architecture discussion.

\textbf{MOCHA} \cite{tanjim2026mocha} observes that agent skills are
``multi-field artifacts subject to hard platform constraints:
description fields are truncated for routing, instruction bodies are
compacted via progressive disclosure, and co-resident skills compete for
limited context windows.'' Skill optimization is therefore not a
single-loss descent problem but an inherently multi-objective one: a
skill must simultaneously maximize task performance and satisfy platform
limits, and optimizing for task performance alone produces skills that
pass benchmarks but violate deployment constraints. MOCHA introduces
\emph{Chebyshev annealing} to navigate the constraint frontier. In
Section III terms, the complexity penalty \(\Omega(z)\) in the gate
condition is precisely the mechanism MOCHA is operationalizing, but
Section III's scalar \(\Omega\) is a simplification: in production, the
constraint surface is a Pareto frontier, not a single penalty term.
MOCHA's empirical contribution is to show that ignoring this frontier
--- treating skill optimization as unconstrained single-loss descent ---
produces skills that degrade on real deployment infrastructure even when
benchmark scores improve.

\textbf{Ratchet} \cite{zhang2026ratchet} makes an argument that is, in
some respects, orthogonal to the gradient-analog framing and for that
reason worth engaging directly. Ratchet's opening finding --- that
LLM-authored skills deliver \(+0.0\) pp over no-skill baselines while
human-curated ones deliver \(+16.2\) pp --- suggests that the bottleneck
in self-evolving skill libraries is not the optimization algorithm (the
mechanism for proposing and accepting edits) but \emph{lifecycle
management}: whether skills are written, retrieved, curated, and retired
correctly. Ratchet introduces four hygiene mechanisms --- outcome-driven
retirement, a bounded active-cap, meta-skill authoring, and pattern
canonicalization --- and shows that these lift the \(+0.0\) pp floor
substantially before any gradient-analog optimizer is applied.

This finding is not a refutation of the Section III picture but a
clarification of the precondition for it. Artifact-layer descent, as
formalized in Section III, presupposes that the scaffold is maintained
well enough for the gate \(G_t\) to have a meaningful signal: if skill
retrieval is so noisy that the wrong skill is injected into every
rollout, the proposal process \(Q_t\) receives corrupted evidence and
the descent guarantee degrades to random search. Ratchet's hygiene
mechanisms are, in Section III vocabulary, prerequisites for operation 1
(residual capture) and operation 2 (mode discovery) to function:
outcome-driven retirement prevents stale artifacts from polluting the
residual signal; the bounded active-cap prevents scaffold bloat from
making retrieval unreliable (the Section X failure mode); pattern
canonicalization collapses duplicate and near-duplicate skills so the
proposal process conditions on a clean, deduplicated artifact set rather
than competing variants of the same pattern. The practical implication
is that the deployment path to Stage-3 artifact-layer descent runs
through Stage-1 hygiene as a \emph{prerequisite}, not as an alternative.
A system with strong hygiene and no optimizer sits at Stage 1 with a
clean scaffold; a system with a strong optimizer and poor hygiene sits
at Stage 2 with a noisy loss signal. Stage 3 requires both.

\textbf{SkillsVote} \cite{liu2026skillsvote} carries the lifecycle
argument to the scale of an open ecosystem. It profiles a million-scale
open-source skill corpus for environment requirements, quality, and
verifiability, synthesizes tasks for the verifiable subset, exposes
skills through agentic library search before execution, and after
execution decomposes each trajectory into skill-linked subtasks whose
outcomes are attributed to skill guidance, agent exploration,
environment, or result signals; only successful, reusable discoveries
pass into evidence-gated updates. Two Section III points follow. First,
attribution is performed at the admission gate rather than inside the
optimizer, which is credit assignment (E.1 op 2) used as a filter on
residual capture (op 1). Second, the redundancy and
environment-sensitivity SkillsVote measures in open skill ecosystems are
the same bloat and drift that Section X names, arising this time from
many authors rather than one loop.

\subsection{S5.8 Cross-task transfer and evidence
radius}\label{s5.8-cross-task-transfer-and-evidence-radius}

Several papers independently measure what Section III calls the
\emph{evidence radius} of a scaffold artifact --- how far an edit
validated in one context continues to improve performance when
redeployed elsewhere --- without using that vocabulary.

\textbf{Training Language Agents to Learn from Experience} {[}ICT;
\cite{shalev2026ict}{]} introduces the In-context Training framework
specifically to evaluate \emph{cross-task self-improvement}: whether
experience distilled from one set of tasks transfers to unseen tasks.
Their finding that ``current reflection-based methods can only
self-correct within a single task instance'' is a measurement of
evidence radius at its lower bound --- a radius of exactly one task ---
and their ICT framework is a systematic attempt to expand that radius
through supervised distillation of reflections into reusable lessons. In
Section III terms, ICT is an attempt to instrument Loop 2 empirically:
to measure whether a locally-validated edit (one that passed the gate on
one task's \(L_D\)) continues to improve performance when the promotion
radius is expanded to cover a distribution of unseen tasks.

\textbf{Combee} \cite{li2026combee} approaches the same measurement
from the scaling direction. ACE \cite{ace2025agentic} established that
context-as-living-artifact (a system prompt that accumulates, refines,
and organizes strategies through agentic interaction) could improve over
static prompts on single-agent tasks. Combee extends this to
high-parallelism settings, arguing that ``existing methods primarily
focus on single-agent or low-parallelism settings,'' which
``fundamentally limits their ability to efficiently learn from a large
set of collected agentic traces.'' The parallel scaling Combee addresses
is a form of Loop-2 scope expansion: when many agents run
simultaneously, the trace pool available for the next proposal step is
drawn from a wider deployment distribution, and the promoted artifact is
being tested against a larger averaging set of patch losses. This is not
cross-organizational Loop-2 promotion, but it is a step toward it, and
Combee's empirical results measure how much of the evidence-radius
benefit can be captured by intra-organization parallelism before the
cross-organizational barrier is reached.

\textbf{CORAL} \cite{qu2026coral} addresses the Loop-2 problem from the
multi-agent open-ended discovery angle. CORAL replaces rigid control
with long-running agents that explore, reflect, and collaborate through
shared persistent memory, asynchronous multi-agent execution, and
heartbeat-based interventions. The shared persistent memory is the
cross-agent promoted artifact: a finding validated in one agent's
exploration that is written to the shared store becomes available to all
agents, and the heartbeat mechanism provides a form of auto-gated Loop-2
promotion. CORAL does not implement the versioned lineage, RBAC, or
rollback that Section IX requires, but it provides the only fully
asynchronous multi-agent Loop-2 promotion mechanism in the current
survey corpus, and its open-ended discovery framing makes the
evidence-radius question concrete.

\textbf{PACE} \cite{ling2026pace} closes this family with a
production-oriented two-timescale design for small language model
agents. The fast timescale adapts prompts, parsers, validators, and
control logic from failure signals within a deployment; the slow
timescale consolidates across episodes. The two-timescale structure maps
cleanly onto the Loop-1 / Loop-2 distinction of Section III: Loop 1
fires at the fast timescale, updating a single deployment's scaffold;
Loop 2 fires at the slow timescale, consolidating across the
deployment's own episodes rather than across deployments. PACE's
contribution is to show that the two-loop architecture is implementable
under the resource constraints of small, frozen production models ---
not only under frontier-model assumptions --- and that the slow/fast
consolidation structure persists even when the optimizer model and the
target model are the same frozen small model.

\textbf{SkillClaw} \cite{ma2026skillclaw} and \textbf{FederatedSkill}
\cite{yang2026federatedskill} are the two systems in the corpus whose
Loop 2 crosses users rather than tasks or agents. SkillClaw treats
cross-user, over-time interaction as the primary improvement signal:
trajectories from many users are aggregated, an autonomous evolver
identifies recurring behavioral patterns and translates them into
skill-set updates, and the updated skills live in a shared repository
synchronized across users; its evaluation protocol runs a day-night
cycle in which simulated concurrent users generate trajectories by day
and the system reflects, proposes bundles, validates, and deploys by
night. FederatedSkill starts from the observation that isolated
single-user task streams lack the diversity to build comprehensive
skills, while raw trajectory sharing compromises privacy and a uniform
global library ignores client heterogeneity. Its unit of communication
is the semantic skill diff (a structured patch over a local library,
shipped with the task reward that produced it); a server-side evolution
agent aggregates patches into a per-client capability matrix (covered,
absorbing, broken, gap) and emits a personalized update for each client,
applying, adapting, or rejecting peer patches with an auditable decision
log that records source patch, reward, and justification for every
modified path. Across 20 task families it reports up to +44.4\% success
and −37.5\% compute against isolated self-evolution. In Section III
terms FederatedSkill is the first surveyed system in which Loop 2 is a
federated aggregation of artifact updates rather than a promotion of a
single artifact, and its decision log is the first lineage record at the
aggregator. It remains inside one operator's federation (clients are
backbone-plus-harness profiles, not organizations), its acceptance
decision uses the source client's reward rather than a held-out
evaluation at the receiving client, and it reports neither RBAC nor
rollback; Section IX's properties (b), (c), and (d) are each partially
or wholly open.

\subsection{S5.9 Synthesis: what the convergence
establishes}\label{s5.9-synthesis-what-the-convergence-establishes}

The four claims stated in Section VIII are supported as follows:

\begin{enumerate}
\def\labelenumi{\arabic{enumi}.}
\item
  \textbf{The scaffold-versus-weights decomposition is independently
  confirmed.} GEPA, Meta-Harness, MOSS, RCL, and PACE all begin from the
  observation that model weights and the external artifact layer are
  distinct adaptation surfaces with different cost, reversibility, and
  scope properties. No surveyed paper disputes this decomposition; the
  variation is in which substrate each paper targets (skill document,
  harness code, system prompt, shared memory).
\item
  \textbf{Credit assignment is the central engineering bottleneck.} AHE,
  CANTANTE, Meta-Harness, \cite{ong2026harnesseval}, HARNESSEVO, and
  SkillsVote all identify credit assignment --- attributing a
  system-level failure to a specific scaffold coordinate --- as the
  problem that distinguishes disciplined Loop-1 evolution from
  trial-and-error. This is the credit-assignment step of Appendix E.1
  (operation 2; E.6), and it remains unsolved in the general
  multi-substrate case.
\item
  \textbf{Lifecycle hygiene is a precondition for descent, not an
  alternative to it.} Ratchet establishes that the gradient-analog
  mechanism of Section III requires a maintained scaffold as a
  precondition. A dirty artifact pool (stale, bloated, miscurated)
  prevents the gate \(G_t\) from receiving a meaningful signal. Hygiene
  and optimization are sequential dependencies, not competing paradigms.
\item
  \textbf{No surveyed system combines Loop-1 automation with
  cross-organizational Loop-2 governance.} SkillOpt has the strongest
  Loop-1 closure and the most systematic evidence-radius measurement,
  but it does not implement cross-organizational promotion, versioned
  lineage with RBAC, or rollback. CORAL implements asynchronous
  cross-agent promotion within one system but not across organizations.
  PACE implements two-timescale consolidation within one deployment.
  Meta-Harness implements Loop-1 closure over the harness substrate but
  no Loop-2 promotion mechanism. FederatedSkill federates skill patches
  across clients under one operator, with a decision log but no RBAC or
  rollback. The four-property gap named in Section IX is narrowed on
  property (b) by FederatedSkill and closed on none; it is, if anything,
  confirmed more strongly by the independent convergence of multiple
  research threads on the same single-organization ceiling.
\end{enumerate}

\subsection{S5.10 SkillOpt against Section IX's four
properties}\label{s5.10-skillopt-against-section-ixs-four-properties}

Mapped onto Section IX's four properties of governed cross-tenant
scaffold optimization:

\begin{itemize}
\tightlist
\item
  \begin{enumerate}
  \def\labelenumi{(\alph{enumi})}
  \tightlist
  \item
    \textbf{Auto-gated Loop 1.} \emph{Satisfied.} The held-out
    selection-split gate at strictly-greater margin instantiates
    \(G_t\).
  \end{enumerate}
\item
  \begin{enumerate}
  \def\labelenumi{(\alph{enumi})}
  \setcounter{enumi}{1}
  \tightlist
  \item
    \textbf{Multi-tenant cross-org promotion.} \emph{Absent.} SkillOpt
    is a single-organization research system; there is no governed
    cross-tenant aggregator.
  \end{enumerate}
\item
  \begin{enumerate}
  \def\labelenumi{(\alph{enumi})}
  \setcounter{enumi}{2}
  \tightlist
  \item
    \textbf{Versioned lineage with RBAC.} \emph{Partial.} The exported
    \texttt{best\_skill.md} is a versioned artifact, and the protected
    slow-update field separates fast and slow consolidation. There is no
    documented RBAC layer or audit-trail policy, and
    \texttt{edit\_apply\_report.json} is an operational log rather than
    an access-controlled lineage system.
  \end{enumerate}
\item
  \begin{enumerate}
  \def\labelenumi{(\alph{enumi})}
  \setcounter{enumi}{3}
  \tightlist
  \item
    \textbf{Rollback.} \emph{Absent.} The training loop selects the best
    validation-gated skill and exports it; no rollback envelope on
    deployed artifacts is reported.
  \end{enumerate}
\end{itemize}

The cross-model, cross-harness, and cross-benchmark transfer experiments
reported by SkillOpt are \emph{empirical evidence-radius measurements}:
they tell us that a skill optimized in one (model, benchmark, harness)
cell continues to improve performance when redeployed in nearby cells
without further optimization. In Section III vocabulary, this is a
measured Loop-2 \emph{radius} for the artifact, not a \emph{governance
mechanism} for promoting the artifact. Both readings are useful:
SkillOpt is the strongest current Loop-1 proof and supplies the first
systematic transfer evidence at the scaffold layer, while the Section IX
cross-tenant governance triad ((b), (c), (d)) remains unfilled.

\subsection{S5.11 Cost and bloat}\label{s5.11-cost-and-bloat}

SkillOpt reports 0.6M to 46.4M training tokens per absolute test-point
gain (their Table 6), with the procedural benchmarks (SpreadsheetBench,
OfficeQA, LiveMath) at the cheap end and search/visual benchmarks
(SearchQA, DocVQA) an order of magnitude more expensive per point. The
deployed artifact remains 379--1,995 tokens regardless. This is, to our
knowledge, the first published cost-per-point figure for scaffold-layer
descent, and it complements the Section X scaffold-bloat discussion (and
the Ratchet hygiene-precondition finding walked in G.7): the textual
learning rate \(\eta_t\) with cosine decay keeps the deployed skill
compact even when the optimizer explores many more edits than are
accepted (only 1--4 edits survive into the final skill across the six
benchmarks). The mechanism that prevents the deployed artifact from
bloating is exactly the complexity penalty \(\Omega(z)\) of Section III,
instantiated as a hard cap on accepted edits per step rather than a soft
regularizer.

\begin{center}\rule{0.5\linewidth}{0.5pt}\end{center}

\section{S6. Claims audit and the rubric-reward
comparison}\label{s6.-claims-audit-and-the-rubric-reward-comparison}

Every absolute claim in the arXiv short v3.2
(\texttt{arxiv/frontier\_and\_localhost/short/paper\_frontier\_and\_localhost\_short.md},
line numbers as of 2026-09-20) of the form \emph{no surveyed system /
only X / first / absent / empty cell}, with its verdict after the
October-2026 literature sweep (three parallel sweeps: skills; harness,
context, and LLMOps governance; memory governance, competing surveys,
rubric rewards; 127 systems added, scored from their own method
sections) and the recount in \texttt{scripts/recount.py}. Verdicts:
\textbf{holds} (unchanged), \textbf{sharpened} (true, now stated more
precisely), \textbf{revised} (wording must change), \textbf{dropped}
(not carried into the IEEE version). Evidence names the system(s) that
forced the verdict; arXiv ids are in \texttt{submission/references.bib}.

\textbf{Corpus after the sweep:} 256 unique systems (130 May + 126
October), 204 core, 52 contrast. No October addition scores 5 on any
substrate. Composite two-loop audit: 14 research systems + 3 industry
exemplars (Sierra Agent OS 2.0, Cognition Devin, and now Letta Context
Repositories, whose shared-memory push carries git lineage but no gate)
pass; 7 October systems promote across contexts without versioning or
lineage and enter only under relaxation (a) or not at all (EvolveNet,
Evo2Team, XRepoSkill, ShareMem, RASO, Spotify, Future AGI (T5 only)).

\textbf{Corrections to the May data applied here
(\texttt{scripts/apply\_corrections.py}):} (1) the row ``Governed
Collaborative Memory / SSGM'' conflated Cuadros et al.~(2605.04264) with
SSGM (Lam et al., 2603.11768); split, SSGM scored 2 (conceptual
framework). (2) MAE and Memento-Skills were scored 5 on S5 in the CSV
while the published table said no orchestration system reaches 5; both
close one loop only, so they are re-scored 4 per the uniform rubric. S5
mean for the May corpus moves from 2.68 to 2.58.

\subsection{S6.The §9 corner (a)--(d)}\label{s6.the-9-corner-ad}

{\def\LTcaptype{none} % do not increment counter
\begin{longtable}[]{@{}
  >{\raggedright\arraybackslash}p{(\linewidth - 6\tabcolsep) * \real{0.2500}}
  >{\raggedright\arraybackslash}p{(\linewidth - 6\tabcolsep) * \real{0.2500}}
  >{\raggedright\arraybackslash}p{(\linewidth - 6\tabcolsep) * \real{0.2500}}
  >{\raggedright\arraybackslash}p{(\linewidth - 6\tabcolsep) * \real{0.2500}}@{}}
\toprule\noalign{}
\begin{minipage}[b]{\linewidth}\raggedright
Line
\end{minipage} & \begin{minipage}[b]{\linewidth}\raggedright
Claim
\end{minipage} & \begin{minipage}[b]{\linewidth}\raggedright
Verdict
\end{minipage} & \begin{minipage}[b]{\linewidth}\raggedright
Evidence
\end{minipage} \\
\midrule\noalign{}
\endhead
\bottomrule\noalign{}
\endlastfoot
70, 120, 235, 246, 272, 574, 578, 592, 698 & No public system combines
(a) auto-gated Loop 1, (b) cross-organization promotion, (c) versioned
lineage with RBAC, (d) rollback & \textbf{holds} & Three sweeps, 126
systems: (b) is unclaimed by any evidenced system (EvolveNet's clients
are benchmark partitions, ``organizations'' only in motivation; MemClaw
is single-tenant in its evaluation; Buildrix is an open community hub
with no gate); (c) appears as lineage without RBAC everywhere (EDDOps,
ACM, MemClaw, Letta, Mem0 history) and as RBAC without versioned
promotion in MCP control planes (agentic-community gateway,
TrueFoundry); (d) is demonstrated by Governed Capability Evolution and
EDDOps and claimed by EvolveNet (round rollback) and EvoUndo
(recoverability as admission), none of which has (a)+(b). \\
110 & ``No surveyed system meets the FedAvg limit; FederatedSkill comes
closest'' & \textbf{sharpened} & EvolveNet (2608.04968) is now the
nearest: broadcast, client-local optimization, server-side merge
accepted iff fixed ≥ broken on a held-out slice, rollback to the
previous shared harness. Still selection-and-merge over program edits,
not averaging; cite both. \\
155 & ``(iv) \ldots{} FederatedSkill comes closest \ldots{} narrowed on
one property, closed on none'' & \textbf{revised} & Two closest
approaches now: FederatedSkill (decision log, no gate at receiver) and
EvolveNet (held-out merge gate and round rollback, no lineage or RBAC).
Say ``narrowed on (b) and (d) by different systems, closed on none.'' \\
242 & ``no member of the §3.5 research families provides (d), and only
FederatedSkill provides even a partial (b)'' & \textbf{revised} & (d):
EvolveNet (round rollback), SkillSmith and SkillAudit (edit-level
rollback inside one agent), Growing Harness and WDA (within-optimization
rollback), EvoUndo (recoverability as an admission test). (b): EvolveNet
and RASO partial. New sentence: ``rollback now appears inside single
deployments and inside one federated round; no system rolls back a
promotion that has crossed an organization boundary, because none makes
such a promotion.'' \\
\end{longtable}
}

\subsection{S6.Matrix cells}\label{s6.matrix-cells}

{\def\LTcaptype{none} % do not increment counter
\begin{longtable}[]{@{}
  >{\raggedright\arraybackslash}p{(\linewidth - 6\tabcolsep) * \real{0.2500}}
  >{\raggedright\arraybackslash}p{(\linewidth - 6\tabcolsep) * \real{0.2500}}
  >{\raggedright\arraybackslash}p{(\linewidth - 6\tabcolsep) * \real{0.2500}}
  >{\raggedright\arraybackslash}p{(\linewidth - 6\tabcolsep) * \real{0.2500}}@{}}
\toprule\noalign{}
\begin{minipage}[b]{\linewidth}\raggedright
Line
\end{minipage} & \begin{minipage}[b]{\linewidth}\raggedright
Claim
\end{minipage} & \begin{minipage}[b]{\linewidth}\raggedright
Verdict
\end{minipage} & \begin{minipage}[b]{\linewidth}\raggedright
Evidence
\end{minipage} \\
\midrule\noalign{}
\endhead
\bottomrule\noalign{}
\endlastfoot
227 & ``Three cells are populated; the rest are empty for structural
reasons'' & \textbf{revised} & Among composite-passing systems the three
cells hold; but single-loop and governance-only systems now occupy
{[}Human × Human{]} (MCP registries: Azure API Center, agentic-community
MCP Gateway \& Registry, TrueFoundry; Buildrix; Enoch), {[}Human ×
None{]} (VerificAgent, revocation guard, Davis et al.~governance
conversion, Continual Harness human phase), {[}Auto+Human × None{]}
(SkillChain, Learning on the Job), {[}None × Auto{]} (Governed
Capability Evolution), {[}None × Auto+Human{]} (EDDOps), {[}None ×
Human{]} (ACM). Restate as: three cells hold every system that passes
the composite audit; the human-authored cells are occupied only by
registries and curation workflows with no session-signal loop. \\
229, 549 & ``{[}Human × Auto{]} is empty by engineering logic \ldots{}
no published example exists'' & \textbf{holds} & No system with
human-authored candidates and automated promotion was found; the
closest, MCP Gateway \& Registry, adds an automated security scan to a
human admission, which is {[}Human(+scan) × Human{]}. \\
229, 551 & ``{[}Human × Human{]} is empty in the surveyed corpus'' &
\textbf{revised} & Occupied by MCP registries and a community hub; see
above. The paper's predicted occupant (solo or small-team workflows) is
also now documented empirically by ``Who Maintains Agent Skills?''
(2609.05677): across five public skill repositories every substantive
edit is human-merged, 62\% AI co-authored, 1 deprecation in 254
edits. \\
229 & ``{[}Human × None{]} is structurally improbable'' &
\textbf{revised} & Occupied: VerificAgent (agent-drafted memories, human
fact-check, frozen bank), the revocation guard of Shen et al., and
human-in-the-loop harness phases. Restate: it is the human-curation
cell, single-loop by construction. \\
231 & ``Nearly every auto-gated system fires at n = 1'' &
\textbf{revised} & Batched or statistical acceptance is now common:
SAGE-gate (paired per-item test), Future AGI (paired bootstrap CI),
AutoSaddler (mini-batch failure signals), MGM, ModularRSI (evidence
across tasks), VaG (3-of-3 critics), EvolveNet (fixed ≥ broken on
held-out), VACE, HackProbe, PROCTOR. Say ``common since mid-2026'' and
cite SAGE-gate and VaG. \\
231 & ``no surveyed system has formalized agreement as a k-of-n
acceptance rule'' & \textbf{revised} & VaG requires all three
heterogeneous critics to pass (3-of-3); SkillSpec is consensus-gated.
Keep the p\^{}k-vs-1/k distinction, drop ``no surveyed system''. \\
231 & ``Only AlphaEvolve demonstrates explicit multi-layer depth'' &
\textbf{sharpened} & Still the only system with cascade evaluator plus
diversity archive; HarnessBank, JIT-Agent, MGM keep archives (lineage),
VaG adds tiered subset selection. Say ``AlphaEvolve remains the only
system combining a cascade evaluator with a diversity archive.'' \\
\end{longtable}
}

\subsection{S6.''Only'' and ``first''}\label{s6.only-and-first}

{\def\LTcaptype{none} % do not increment counter
\begin{longtable}[]{@{}
  >{\raggedright\arraybackslash}p{(\linewidth - 6\tabcolsep) * \real{0.2500}}
  >{\raggedright\arraybackslash}p{(\linewidth - 6\tabcolsep) * \real{0.2500}}
  >{\raggedright\arraybackslash}p{(\linewidth - 6\tabcolsep) * \real{0.2500}}
  >{\raggedright\arraybackslash}p{(\linewidth - 6\tabcolsep) * \real{0.2500}}@{}}
\toprule\noalign{}
\begin{minipage}[b]{\linewidth}\raggedright
Line
\end{minipage} & \begin{minipage}[b]{\linewidth}\raggedright
Claim
\end{minipage} & \begin{minipage}[b]{\linewidth}\raggedright
Verdict
\end{minipage} & \begin{minipage}[b]{\linewidth}\raggedright
Evidence
\end{minipage} \\
\midrule\noalign{}
\endhead
\bottomrule\noalign{}
\endlastfoot
145, 299, 564, 682 & CORAL ``the only asynchronous multi-agent Loop-2
mechanism in the corpus'' & \textbf{revised} & Evo2Team (2609.34135)
promotes routing and communication skills across 4--32-agent teams with
target-side confirmation; ShareMem pools memory across users. Say
``CORAL was the only one in the May corpus; Evo2Team and ShareMem add
cross-team and cross-user promotion without versioning.'' \\
213 & AlphaEvolve Strassen claim & \textbf{holds} (attributed) &
Third-party; attribute to the AlphaEvolve report in the same
sentence. \\
215 & ``CASCADE reports the largest within-benchmark ablation gain in
the corpus'' & \textbf{dropped} & Not re-checked across 127 new systems;
comparing ablation gains across benchmarks is not meaningful. Removed
from the IEEE text. \\
254, 317 & ChronoMem and rollback repair ``the first surveyed instances
of the rollback envelope'' & \textbf{revised} & Rollback of memory is
now in VerMem (soft-delete/restore ops), Mi-Memory (design), Letta (git
history); and shown unsafe in checkpoint-rollback studies (2608.29381,
2608.15939). Say ``among the first'' and cite the counter-evidence. \\
316 & Cuadros ``only surveyed memory system with full provenance +
versioning + correction'' & \textbf{holds} & Memory governance additions
(TOKI, PPMF, TMA-NM, MemClaw, He \& Yu) each carry provenance or lineage
without correction loops or versioning, and the auto-gated memory loops
(Recuris, MemGuard, VerMem, Compiled Memory) carry no lineage. \\
610, 630 & SkillOpt ``cleanest published instance'' / ``first to
instantiate all four simultaneously'' & \textbf{sharpened} &
SkillOpt-Lite (2607.03451) removes two of the four and matches it;
VALVE, SkillSmith, SkillAudit instantiate held-out gates with rollback.
Say ``the first system we found that instantiates all four, and its own
follow-up shows two of them are stabilizers rather than
necessities.'' \\
686 & FederatedSkill ``first \ldots{} federated aggregation \ldots{}
first lineage record at the aggregator'' & \textbf{holds} & EvolveNet
(Aug 2026) postdates it and has no lineage record. \\
713 & SkillOpt ``first published cost-per-point figure'' &
\textbf{holds} & No earlier figure found. \\
\end{longtable}
}

\subsection{S6.Absences}\label{s6.absences}

{\def\LTcaptype{none} % do not increment counter
\begin{longtable}[]{@{}
  >{\raggedright\arraybackslash}p{(\linewidth - 6\tabcolsep) * \real{0.2500}}
  >{\raggedright\arraybackslash}p{(\linewidth - 6\tabcolsep) * \real{0.2500}}
  >{\raggedright\arraybackslash}p{(\linewidth - 6\tabcolsep) * \real{0.2500}}
  >{\raggedright\arraybackslash}p{(\linewidth - 6\tabcolsep) * \real{0.2500}}@{}}
\toprule\noalign{}
\begin{minipage}[b]{\linewidth}\raggedright
Line
\end{minipage} & \begin{minipage}[b]{\linewidth}\raggedright
Claim
\end{minipage} & \begin{minipage}[b]{\linewidth}\raggedright
Verdict
\end{minipage} & \begin{minipage}[b]{\linewidth}\raggedright
Evidence
\end{minipage} \\
\midrule\noalign{}
\endhead
\bottomrule\noalign{}
\endlastfoot
62 & ``production deployments sit at Stage 1 or 2; none reaches Stage
3'' & \textbf{sharpened} & Spotify (2609.30297) runs automated prompt
and planning optimization gated by human launch and A/B metrics;
SkillChain (Alibaba) runs judge-monotone edits behind a human reflection
gate with reported agreement. Both are Stage 3 on Loop 1 with a human
Loop 2; say ``production deployments reach Stage 3 on the local loop
only behind a human promotion gate''. \\
188 & S1 score-5 none; ``no instruction system closes both loops'' &
\textbf{holds} & S1 n=15, max 4; adoption studies cut the other way
(Gloaguen: context files raise cost \textgreater20\% without raising
success; Galster: context files dominate configuration, skills rare). \\
189 & ``Anthropic's skills repository is one-pool with PR review but no
automated eval gate'' & \textbf{holds} & Confirmed by 2609.05677
(human-merged, no eval gate observed). \\
190 & ``correctable, versioned memory is rare; production memory lacks
failure-triggered correction and versioned lineage'' & \textbf{revised}
& Letta Context Repositories (git-backed memory, every edit committed,
shared memory auto-pushed), Mem0 per-memory history and feedback
endpoint, Zep bitemporal invalidation: lineage is now common in
production memory; failure-triggered correction is not, and 2609.08258
shows revoked facts still outrank their replacements in five backends.
Say ``versioned lineage has arrived in production memory; enforced
correction has not.'' \\
192 & S5 score-5 none; ``topology evolution is rare'' &
\textbf{sharpened} & After the correction S5 has no 5 in either corpus.
Gated topology evolution now exists: TOCOMAS retains only candidates
satisfying structural invariants; Evo2Team promotes across teams;
OptiMAS, AgentGrad, HarnessForge update S5 under evaluation. Say ``rare
and, where it exists, gated by evaluation alone except TOCOMAS.'' \\
196 & ``no system matures all six at once'' & \textbf{holds} & No
October addition scores 5 on any substrate. \\
198a & ``Failure-triggered memory is rare in production'' &
\textbf{holds} & One field report (SIx Harness, N=1, self-reported);
Letta dreaming writes from conversation history, not failures. \\
198b & ``Cross-customer skill promotion is nearly absent'' &
\textbf{holds} & ShareMem, EvolveNet, XRepoSkill, RASO promote across
users, clients, or repositories inside one operator; none across
customers or organizations with evidence. \\
198c & ``Memory versioning is mostly unsolved outside Cuadros et al.~and
ChronoMem'' & \textbf{revised} & See 190: Letta, Mem0 history, Zep
bitemporal facts, TOKI audit rows, MemClaw supersession chains. \\
198d & ``Scaffold-level curriculum is unspecified in nearly every
surveyed system'' & \textbf{sharpened} & Task-CoEvolve
(variance-weighted task sampling) and AutoSaddler (mini-batch failure
signals) specify it; still the exception. \\
198e & ``Held-out evals at Loop-2 promotion are absent'' &
\textbf{revised} & Receiving-side evaluation now exists: ShareMem
(evaluation tasks belong to a different user), XRepoSkill (independent
verification in ≥2 repositories), Evo2Team (target-side confirmation; 4
of 22 transfers still fail on new graphs), VACE (candidate evaluated
with the receiving model fixed). EvolveNet explicitly has none;
Federated Agent Optimization (2610.01195) proposes a receiving-client
gate without implementing it. \\
206, 213 & research full-auto: ``enterprise governance is uniformly
absent'' & \textbf{holds} & No October {[}Auto × Auto{]} system reports
RBAC or audited promotion. \\
208, 217 & governance-first: ``no patch-local inner loop'' &
\textbf{holds} & EDDOps, ACM, GCE, MCP control planes all have Loop 1 =
None. \\
219 & ``AGENTS.md in over 60,000 repositories'' & \textbf{revised} &
Newer counts: 2,853-repository study (Galster), 128,018-project adoption
study (Robbes: 22--29\% of coding-agent projects), 36,710 repositories
mined for loop engineering (Lulla 2608.21884). Replace the single count
with these studies and the cost finding (Gloaguen). \\
264 & ``no surveyed system treats pruning or expiration as a first-class
gated operation (outside Ratchet, SkillOpt, ChronoMem)'' &
\textbf{revised} & Retirement is a gated operation in at least eleven
skill systems (Library Drift: retire when n ≥ 100 and contribution ≤
−0.10 with a cap of 50; SkillSmith; SkillOps; ASSAY; HiSME; SkeMex;
MUSE-Autoskill; SkillMAS; SafeEvolve; Double Ratchet; SkillGen) and
expiry is a product control in Mem0, Zep, and LangGraph Platform. Say
``pruning and expiry became first-class in 2026; what remains rare is
pruning of a promoted artifact's descendants, which VaG shows is
structurally irreversible.'' \\
266 & failure surface ``largely unmitigated in production'' &
\textbf{holds} & Revocation unenforced in five production backends
(2609.08258); checkpoint rollback unsafe in three frameworks
(2608.29381); no evidentiary record in 16 deployed frameworks (Hearsay
2609.32495). \\
375 & multi-sample agreement ``under-explored; SkillOpt-Lite nearest'' &
\textbf{revised} & See 231: SAGE-gate, VaG, SkillSpec. \\
692 & ``No surveyed paper disputes the scaffold-vs-weights
decomposition'' & \textbf{holds} & SIA (2605.27276) updates harness and
weights together but keeps them distinct levers. \\
694 & credit assignment ``remains unsolved in the general
multi-substrate case'' & \textbf{holds} & HARNESSEVO (slot attribution),
AgentGrad (intervention-based), CANTANTE (contrastive) each solve one
substrate's case. \\
\end{longtable}
}

\subsection{S6.Section IV claims (checked against the rescored
corpus)}\label{s6.section-iv-claims-checked-against-the-rescored-corpus}

{\def\LTcaptype{none} % do not increment counter
\begin{longtable}[]{@{}
  >{\raggedright\arraybackslash}p{(\linewidth - 6\tabcolsep) * \real{0.2500}}
  >{\raggedright\arraybackslash}p{(\linewidth - 6\tabcolsep) * \real{0.2500}}
  >{\raggedright\arraybackslash}p{(\linewidth - 6\tabcolsep) * \real{0.2500}}
  >{\raggedright\arraybackslash}p{(\linewidth - 6\tabcolsep) * \real{0.2500}}@{}}
\toprule\noalign{}
\begin{minipage}[b]{\linewidth}\raggedright
Claim
\end{minipage} & \begin{minipage}[b]{\linewidth}\raggedright
Statement
\end{minipage} & \begin{minipage}[b]{\linewidth}\raggedright
Verdict
\end{minipage} & \begin{minipage}[b]{\linewidth}\raggedright
Evidence
\end{minipage} \\
\midrule\noalign{}
\endhead
\bottomrule\noalign{}
\endlastfoot
IV-1 & Every {[}Auto × Auto{]} system with a per-edit estimator gate
runs on a C1 signal or a simulated C2 environment; {[}Auto × Auto{]}
systems in C4 domains gate by proposal pruning & \textbf{holds} & May
{[}Auto × Auto{]}: DGM, Voyager, AlphaEvolve, SkillOpt, SkillRL, SAGE,
EvolveR, CoEvoSkills, AutoAgent, AutoManual (C1/simulated C2);
NanoResearch (semantic merge, C4) and CORAL (heartbeat, C4) prune.
October: EvolveNet, Evo2Team, XRepoSkill (C1 code/benchmark signals);
ShareMem (C2 Mind2Web); RASO (target execution feedback, C1/C2); Letta
(no gate). \\
IV-2 & {[}Auto × Human{]} is populated only by C2/C3 systems &
\textbf{holds} & SkillForge (C3 judge, human Loop 2), CASCADE (C3),
Sierra (C2/C3), Devin (C2 merge outcome), Spotify (C2 A/B behind human
launch). No C1 system occupies the cell. \\
IV-3 & Rubric-as-reward builds C3 gates on the weight side;
scaffold-side judge calibration under the reflexive condition is rare &
\textbf{sharpened} & 14 reward papers: 8 fix the rubric or judge during
policy optimization, 5 co-evolve it (RLPR by construction, EvoRubrics,
SCORE, RubricEM, Rubric-ARM at the reward-model stage); only 5 report
any human-agreement number and none of the co-evolving ones does.
Scaffold side: SkillForge (\textgreater90\% agreement), SkillChain (κ
0.72--0.91), SkillCoach (trajectory-filter consistency 82→96\% against
human gold) now report calibration; none calibrates a judge that is
itself being edited. \\
\end{longtable}
}

\subsection{S6.Recount output (scripts/recount.py, merged corpus, all
rows)}\label{s6.recount-output-scriptsrecount.py-merged-corpus-all-rows}

\begin{verbatim}
rows: 269  unique systems: 256  core: 204  contrast: 52
substrate            n   mean  score-5 systems
INTEGRATED           52   3.81  AutoAgent, NanoResearch, SkillRL
S1_Instructions      15   3.07  (none)
S2_Skills            50   3.34  COSPLAY, SAGE
S3_Memory            50   3.06  Governed Collaborative Memory (Cuadros et al.)
S4_Tools             20   3.65  Harvey, Hippocratic AI, MCP
S5_Orchestration     37   3.08  (none)
S6_Governance        45   3.02  AGENTS.md / AAIF, Braintrust, LangSmith Hub, Vellum
May corpus alone (after corrections): 130 / 99 / 31; S5 mean 2.58.
\end{verbatim}

\bibliography{references.bib}

\bibliographystyle{IEEEtran}\bibliography{references}
\end{document}
